% Chapter 3, Section 5 _Linear Algebra_ Jim Hefferon
%  http://joshua.smcvt.edu/linearalgebra
%  2001-Jun-12
\section{Perubahan Basis}
\index{perubahan basis|(}
Representasi berubah mengikuti basisnya.
Sebagai contoh, terhadap basis $\stdbasis_2$ dan
\begin{equation*}
  B=\sequence{\colvec{1 \\ 1},\colvec{1 \\ -1}}
\end{equation*}
$\vec{e}_1\in\Re^2$ memiliki representasi-representasi berbeda berikut.
\begin{equation*}
  \rep{\vec{e}_1}{\stdbasis_2}=\colvec{1 \\ 0}
  \qquad
  \rep{\vec{e}_1}{B}=\colvec{1/2 \\ 1/2}
\end{equation*}
Hal yang sama berlaku bagi pemetaan:
terhadap pasangan basis $\stdbasis_2,\stdbasis_2$ dan $\stdbasis_2,B$,
pemetaan identitas memiliki representasi berikut.
\begin{equation*}
  \rep{\text{id}}{\stdbasis_2,\stdbasis_2}=
   \begin{mat}[r]
     1  &0  \\
     0  &1
   \end{mat}
   \qquad
  \rep{\text{id}}{\stdbasis_2,B}=
   \begin{mat}[r]
     1/2  &1/2  \\
     1/2  &-1/2
   \end{mat}
\end{equation*}
Bagian ini menunjukkan cara menerjemahkan di antara representasi-representasi.
Artinya, kita akan menghitung bagaimana representasi berubah ketika basisnya
berubah.












\subsection{Mengubah Representasi Vektor}
Ketika mengubah
$\rep{\vec{v}}{B}$ menjadi $\rep{\vec{v}}{D}$
vektor yang mendasarinya, $\vec{v}$, tidak berubah.
Jadi, penerjemahan antara kedua cara menyatakan vektor ini dilakukan oleh
pemetaan identitas pada ruang tersebut, yang dideskripsikan sedemikian sehingga
vektor-vektor ruang domain direpresentasikan terhadap $B$ dan vektor-vektor
ruang kodomain direpresentasikan terhadap $D$.
\index{diagram panah}
%<*ChangeRepresentationOfVectorArrowDiagram>
\begin{equation*}
  \begin{CD}
    V_{\wrt{B}}                      \\
    @V{\text{\scriptsize$\identity$}} VV   \\
    V_{\wrt{D}}
  \end{CD}
\end{equation*}
(Diagram ini vertikal agar sesuai dengan diagram-diagram dalam subbagian
berikutnya.)
%</ChangeRepresentationOfVectorArrowDiagram>

\begin{definition}  \label{df:ChangeOfBasisMatrix}
%<*df:ChangeOfBasisMatrix>
\definend{Matriks perubahan basis}\index{vektor!representasi}%
\index{matriks!perubahan basis}\index{basis!perubahan}
untuk basis \( B,D\subset V \) adalah representasi pemetaan identitas
\( \map{\identity}{V}{V} \) terhadap kedua basis tersebut.
\begin{equation*}
  \rep{\identity}{B,D}=
  \begin{pmat}{c@{\hspace{1em}}c@{\hspace{1em}}c}
     \vdots                  &             &\vdots                  \\
     \rep{\vec{\beta}_1}{D}  &\;\cdots\;   &\rep{\vec{\beta}_n}{D}  \\
     \vdots                  &             &\vdots
  \end{pmat}
\end{equation*}
%\( B=\basis{\beta}{n} \).
%</df:ChangeOfBasisMatrix>
\end{definition}

\begin{remark}
Nama yang lebih tepat adalah `matriks perubahan representasi', tetapi nama di
atas merupakan nama standar.
\end{remark}

Hasil berikut mendukung definisi tersebut.

\begin{lemma}  \label{le:ChBasisMatDoesChBases}
%<*lm:ChBasisMatDoesChBases>
Untuk mengubah representasi vektor~$\vec{v}$ terhadap \( B \) menjadi
representasinya terhadap \( D \), gunakan matriks perubahan basis.
\begin{equation*}
  \rep{\identity}{B,D}\,\rep{\vec{v}}{B}=\rep{\vec{v}}{D}
\end{equation*}
Sebaliknya, jika perkalian kiri dengan suatu matriks mengubah basis,
$M\cdot\rep{\vec{v}}{B}=\rep{\vec{v}}{D}$
maka $M$ merupakan matriks perubahan basis.
%</lm:ChBasisMatDoesChBases>
\end{lemma}

\begin{proof}
%<*pf:ChBasisMatDoesChBases>
Kalimat pertama berlaku karena perkalian matriks-vektor merepresentasikan
penerapan pemetaan, sehingga
\( \rep{\identity}{B,D}\cdot\rep{\vec{v}}{B}=\rep{\,\identity(\vec{v})\,}{D}
  =\rep{\vec{v}}{D} \) bagi setiap \( \vec{v} \).
Untuk kalimat kedua, terhadap $B,D$, matriks $M$ merepresentasikan pemetaan
linear yang tindakannya memetakan setiap vektor ke dirinya sendiri, sehingga
merupakan pemetaan identitas.
%</pf:ChBasisMatDoesChBases>
\end{proof}

\begin{example}
Dengan basis-basis berikut bagi \( \Re^2 \),
\begin{equation*}
  B=
  \sequence{
            \colvec[r]{2 \\ 1},
            \colvec[r]{1 \\ 0} }
  \qquad
  D=
  \sequence{
            \colvec[r]{-1 \\ 1},
            \colvec[r]{1 \\ 1} }
\end{equation*}
karena
\begin{equation*}
  \rep{\,\identity(\colvec[r]{2 \\ 1})}{D}
  =\colvec[r]{-1/2 \\ 3/2}_D
  \qquad
  \rep{\,\identity(\colvec[r]{1 \\ 0})}{D}
  =\colvec[r]{-1/2 \\ 1/2}_D
\end{equation*}
matriks perubahan basisnya adalah sebagai berikut.
\begin{equation*}
  \rep{\rm id}{B,D}
  =
    \begin{mat}[r]
       -1/2  &-1/2  \\
        3/2  &1/2
    \end{mat}
\end{equation*}
Sebagai contoh, berikut adalah representasi $\vec{e}_2$
\begin{equation*}
  \rep{\colvec[r]{0 \\ 1} }{B}
  =\colvec[r]{1 \\ -2}
  % \qquad
  % \rep{\colvec[r]{0 \\ 1} }{D}
  % =\colvec[r]{1/2 \\ 1/2}
\end{equation*}
dan matriks tersebut melakukan pengubahannya.
\begin{equation*}
    \begin{mat}[r]
       -1/2  &-1/2  \\
        3/2  &1/2
    \end{mat}
  \colvec[r]{1 \\ -2}
  =
  \colvec[r]{1/2 \\ 1/2}
\end{equation*}
Pemeriksaan bahwa vektor di sebelah kanan adalah $\rep{\vec{e}_2 }{D}$
mudah dilakukan.
\end{example}

Kita menutup subbagian ini dengan mengenali matriks-matriks perubahan basis
sebagai suatu himpunan yang telah dikenal.

\begin{lemma}    \label{le:NonSingIsChBasis}
%<*lm:NonSingIsChBasis>
Suatu matriks mengubah basis jika dan hanya jika matriks tersebut nonsingular.
%</lm:NonSingIsChBasis>
\end{lemma}

\begin{proof}
%<*pf:NonSingIsChBasis0>
Untuk arah `hanya jika', jika perkalian kiri dengan suatu matriks mengubah
basis, maka matriks tersebut merepresentasikan fungsi yang dapat diinvers,
karena kita dapat menginvers fungsi itu dengan mengubah basisnya kembali.
Karena merepresentasikan fungsi yang dapat diinvers, matriks itu sendiri dapat
diinvers sehingga nonsingular.
%</pf:NonSingIsChBasis0>

%<*pf:NonSingIsChBasis1>
Untuk arah `jika', kita akan menunjukkan bahwa sebarang matriks nonsingular
$M$ melakukan operasi perubahan basis dari sebarang basis awal $B$
(yang memiliki~$n$ vektor, dengan matriks berukuran $\nbyn{n}$) ke suatu
basis akhir.

% The matrix is nonsingular so it will Gauss-Jordan reduce to the
% identity.
Jika matriksnya adalah identitas~$I$, pernyataan tersebut jelas.
Jika tidak, karena matriksnya nonsingular,
Korolari~IV.\ref{cor:ReducViaMatrices} menyatakan bahwa terdapat matriks-matriks
reduksi elementer sedemikian sehingga $R_r\cdots R_1\cdot M=I$ dengan
$r\geq 1$.
Matriks elementer dapat diinvers dan inversnya juga elementer, sehingga
mengalikan kedua ruas persamaan tersebut dari kiri dengan ${R_r}^{-1}$,
lalu ${R_{r-1}}^{-1}$, dan seterusnya, menyatakan $M$ sebagai hasil kali
matriks-matriks elementer $M={R_1}^{-1}\cdots {R_r}^{-1}$.
% (We've combined ${R_r}^{-1}I$ to make ${R_r}^{-1}$; 
% because $r\geq 1$ we can always have the $I$ disappear in this way.)
%</pf:NonSingIsChBasis1>

%<*pf:NonSingIsChBasis2>
Pembuktian selesai jika kita menunjukkan bahwa matriks-matriks elementer
mengubah suatu basis menjadi basis lain, sebab kemudian ${R_r}^{-1}$ mengubah
$B$ menjadi suatu basis lain $B_r$, ${R_{r-1}}^{-1}$ mengubah $B_r$ menjadi
suatu $B_{r-1}$, dan seterusnya.
Kita akan membahas ketiga jenis matriks elementer secara terpisah; ingat kembali
notasi untuk ketiganya.
%</pf:NonSingIsChBasis2>
% This equation doesn't fit well on slide, so I copied it over an tweaked it.
\begin{equation*}
  M_{i}(k)
    \colvec{
       c_1     \\
       \vdots  \\
       c_i    \\
       \vdots  \\
       c_n  }
  =
    \colvec{
       c_1     \\
       \vdots  \\
       kc_i    \\
       \vdots  \\
       c_n  }
  \quad
  P_{i,j}
    \colvec{
       c_1     \\
       \vdots  \\
       c_i    \\
       \vdots  \\
       c_j    \\
       \vdots  \\
       c_n  }
  =
    \colvec{
       c_1     \\
       \vdots  \\
       c_j    \\
       \vdots  \\
       c_i    \\
       \vdots  \\
       c_n  }
  \quad
  C_{i,j}(k)
    \colvec{
       c_1     \\
       \vdots  \\
       c_i    \\
       \vdots  \\
       c_j    \\
       \vdots  \\
       c_n  }
  =
    \colvec{
       c_1     \\
       \vdots  \\
       c_i    \\
       \vdots  \\
       kc_i+c_j    \\
       \vdots  \\
       c_n  }
\end{equation*}
%<*pf:NonSingIsChBasis3>
Menerapkan matriks perkalian baris~$M_{i}(k)$ mengubah representasi terhadap
\( \sequence{\vec{\beta}_1,\dots,\vec{\beta}_i,\dots,\vec{\beta}_n } \)
menjadi representasi terhadap
\( \sequence{\vec{\beta}_1,\dots,(1/k)\vec{\beta}_i,\dots,\vec{\beta}_n } \).
\begin{multline*}
   \vec{v}= c_1\cdot\vec{\beta}_1+\dots+c_i\cdot\vec{\beta}_i
                                   +\dots+c_n\cdot\vec{\beta}_n  
   \\  \mapsto\;                                                       
    c_1\cdot\vec{\beta}_1+\dots+kc_i\cdot(1/k)\vec{\beta}_i+\dots
                                  +c_n\cdot\vec{\beta}_n=\vec{v}     
\end{multline*}
Barisan kedua merupakan basis karena barisan pertama merupakan basis dan karena
pembatasan $k\neq 0$ dalam definisi matriks perkalian baris.
%</pf:NonSingIsChBasis3>
%<*pf:NonSingIsChBasis4>
Demikian pula, perkalian kiri dengan matriks pertukaran baris $P_{i,j}$
mengubah representasi terhadap basis
\( \sequence{\vec{\beta}_1,\dots,\vec{\beta}_i,\dots,\vec{\beta}_j,
  \dots,\vec{\beta}_n } \)
menjadi representasi terhadap basis berikut
\( \sequence{\vec{\beta}_1,\dots,\vec{\beta}_j,\dots,\vec{\beta}_i,
  \dots,\vec{\beta}_n } \).
\begin{multline*}
   \vec{v}= c_1\cdot\vec{\beta}_1+\dots+c_i\cdot\vec{\beta}_i
                   +\dots+c_j\vec{\beta}_j+\dots+c_n\cdot\vec{\beta}_n  
   \\  \mapsto\;                                                       
    c_1\cdot\vec{\beta}_1+\dots+c_j\cdot\vec{\beta}_j+\dots
               +c_i\cdot\vec{\beta}_i+\dots+c_n\cdot\vec{\beta}_n=\vec{v}     
\end{multline*}
%</pf:NonSingIsChBasis4>
%<*pf:NonSingIsChBasis5>
Selanjutnya, representasi terhadap
\( \sequence{\vec{\beta}_1,\dots,\vec{\beta}_i,\dots,\vec{\beta}_j,
  \dots,\vec{\beta}_n } \)
berubah melalui perkalian kiri dengan matriks kombinasi baris
$C_{i,j}(k)$ menjadi representasi terhadap
\( \sequence{\vec{\beta}_1,\dots,\vec{\beta}_i-k\vec{\beta}_j,
  \dots,\vec{\beta}_j,\dots,\vec{\beta}_n } \)
\begin{multline*}
   \vec{v}= c_1\cdot\vec{\beta}_1+\dots+c_i\cdot\vec{\beta}_i
                       +c_j\vec{\beta}_j+\dots+c_n\cdot\vec{\beta}_n  
   \\  \mapsto\;                                                       
    c_1\cdot\vec{\beta}_1+\dots+c_i\cdot(\vec{\beta}_i-k\vec{\beta}_j)+\dots
          +(kc_i+c_j)\cdot\vec{\beta}_j+\dots+c_n\cdot\vec{\beta}_n=\vec{v} 
\end{multline*}
(definisi $C_{i,j}(k)$ menetapkan bahwa \( i\neq j \) dan
\( k\neq 0 \)).
%</pf:NonSingIsChBasis5>
\end{proof}

\begin{corollary}  \label{co:MatrixNonsingularIffChangesBasis}
%<*co:MatrixNonsingularIffChangesBasis>
Suatu matriks nonsingular jika dan hanya jika matriks tersebut merepresentasikan
pemetaan identitas terhadap suatu pasangan basis.
%</co:MatrixNonsingularIffChangesBasis>
\end{corollary}

% Just as this subsection shows how to translate among representations of
% vectors 
% the next subsection will
% show how to translate among representations of 
% maps, that is, how to change
% $\rep{h}{B,D}$ to $\rep{h}{\hat{B},\hat{D}}$.
% The above corollary is a special case of this, where the domain and range are
% the same space and where the map is the identity map.


\begin{exercises}
  \recommended \item 
    Dalam \( \Re^2 \), dengan
    \begin{equation*}
      D=\sequence{\colvec[r]{2 \\ 1},\colvec[r]{-2 \\ 4}}
    \end{equation*}
    temukan matriks perubahan basis dari \( D \) ke \( \stdbasis_2 \) dan
    dari \( \stdbasis_2 \) ke \( D \).
    Kalikan kedua matriks tersebut.
    \begin{answer}
      Agar matriks mengubah basis dari $D$ ke $\stdbasis_2$, kita memerlukan
      $\rep{\identity(\vec{\delta}_1)}{\stdbasis_2}
        =\rep{\vec{\delta}_1}{\stdbasis_2}$ 
      dan
      $\rep{\identity(\vec{\delta}_2)}{\stdbasis_2}
        =\rep{\vec{\delta}_2}{\stdbasis_2}$.
      Tentu saja, representasi suatu vektor dalam $\Re^2$ terhadap basis
      standar mudah diperoleh.
      \begin{equation*}
        \rep{\vec{\delta}_1}{\stdbasis_2}=\colvec[r]{2 \\ 1}
        \qquad
        \rep{\vec{\delta}_2}{\stdbasis_2}=\colvec[r]{-2 \\ 4}
      \end{equation*}
      Menyusun kedua vektor itu sebagai kolom-kolom matriks perubahan basis
      menghasilkan
      \begin{equation*}
        \rep{\identity}{D,\stdbasis_2}
        =\begin{mat}[r]
          2     &-2    \\
          1     &4
        \end{mat}
      \end{equation*}
      Untuk matriks perubahan basis dalam arah sebaliknya, kita dapat
      menghitung $\rep{\identity(\vec{e}_1)}{D}=\rep{\vec{e}_1}{D}$ dan
      $\rep{\identity(\vec{e}_2)}{D}=\rep{\vec{e}_2}{D}$ (perhitungan ini
      rutin), atau mengambil invers matriks di atas.
      Rumus invers matriks $\nbyn{2}$ membuat perhitungan ini mudah.
      \begin{equation*}
        \rep{\identity}{\stdbasis_2,D}=
        \frac{1}{10}\cdot\begin{mat}[r]
          4  &2  \\
         -1  &2
        \end{mat}
        =\begin{mat}[r]
          4/10  &2/10  \\
         -1/10  &2/10
        \end{mat}
      \end{equation*}
    \end{answer}
 \item Manakah di antara matriks-matriks berikut yang dapat digunakan untuk
   mengubah basis?
   \begin{exparts*}
     \partsitem 
       $\begin{mat}
          1 &2 \\
          3 &4 
        \end{mat}$
     \partsitem 
       $\begin{mat}
          0 &-1  \\
          1 &-1
        \end{mat}$
     \partsitem 
       $\begin{mat}
          2 &3 &-1 \\
          0 &1 &0
        \end{mat}$
     \partsitem 
       $\begin{mat}
          2 &3 &-1 \\
          0 &1 &0  \\
          4 &7 &-2
        \end{mat}$
     \partsitem 
       $\begin{mat}
          0 &2 &0 \\
          0 &0 &6 \\
          1 &0 &0
        \end{mat}$
   \end{exparts*}
   \begin{answer}
     Jika suatu matriks nonsingular, matriks tersebut dapat menjadi matriks
     perubahan basis. Untuk semua matriks ini, kita dapat memeriksanya secara
     langsung.
     \begin{exparts}
       \partsitem Matriks ini nonsingular karena baris kedua bukan kelipatan
          baris pertama.
       \partsitem Matriks ini nonsingular.
       \partsitem Matriks ini singular, sebab suatu matriks harus persegi agar
          dapat bersifat nonsingular.
       \partsitem Matriks ini singular karena dua kali baris pertama ditambah
          baris kedua sama dengan baris ketiga.
       \partsitem Nonsingular.
     \end{exparts}
   \end{answer}
  \recommended \item \label{exer:ChngeBasMatsRTwo}
    Temukan matriks perubahan basis untuk \( B,D\subseteq\Re^2 \).
    \begin{exparts*}
      \partsitem \( B=\stdbasis_2 \),
        \( D=\sequence{\vec{e}_2,\vec{e}_1} \)
      \partsitem \( B=\stdbasis_2 \),
        \( D=\sequence{\colvec[r]{1 \\ 2},\colvec[r]{1 \\ 4}} \)
      \partsitem  \( B=\sequence{\colvec[r]{1 \\ 2},\colvec[r]{1 \\ 4}} \),
        \( D=\stdbasis_2 \)
      \partsitem  \( B=\sequence{\colvec[r]{-1 \\ 1},\colvec[r]{2 \\ 2}} \),
        \( D=\sequence{\colvec[r]{0 \\ 4},\colvec[r]{1 \\ 3}} \)
    \end{exparts*}
    \begin{answer}
      Susun
      $\rep{\identity(\vec{\beta}_1)}{D}=\rep{\vec{\beta}_1}{D}$
      dan $\rep{\identity(\vec{\beta}_2)}{D}=\rep{\vec{\beta}_2}{D}$
      sebagai kolom-kolom untuk membentuk matriks perubahan basis
      $\rep{\identity}{B,D}$.
      \begin{exparts*}
        \partsitem \( \begin{mat}[r]
                   0  &1  \\
                   1  &0
                 \end{mat} \)
        \partsitem \( \begin{mat}[r]
                   2  &-1/2  \\
                   -1 &1/2
                 \end{mat} \)
        \partsitem \( \begin{mat}[r]
                   1  &1     \\
                   2  &4
                 \end{mat} \)
        \partsitem \( \begin{mat}[r]
                   1  &-1    \\
                  -1  &2
                 \end{mat} \)
      \end{exparts*}  
    \end{answer}
  \recommended \item
    Temukan matriks perubahan basis untuk setiap
    \( B,D\subseteq\polyspace_2 \).
    \begin{exparts*}
      \partsitem \( B=\sequence{1,x,x^2},
               D=\sequence{x^2,1,x} \)
      \partsitem \( B=\sequence{1,x,x^2},
               D=\sequence{1,1+x,1+x+x^2} \)
      \partsitem \( B=\sequence{2,2x,x^2},
               D=\sequence{1+x^2,1-x^2,x+x^2} \)
    \end{exparts*}
    \begin{answer}
      Vektor-vektor
      $\rep{\identity(\vec{\beta}_1)}{D}=\rep{\vec{\beta}_1}{D}$,
      $\rep{\identity(\vec{\beta}_2)}{D}=\rep{\vec{\beta}_2}{D}$,
      dan $\rep{\identity(\vec{\beta}_3)}{D}=\rep{\vec{\beta}_3}{D}$
      membentuk kolom-kolom matriks perubahan basis
      $\rep{\identity}{B,D}$.
      \begin{exparts*}
        \partsitem \( \begin{mat}[r]
                   0  &0  &1  \\
                   1  &0  &0  \\
                   0  &1  &0
                 \end{mat} \)
        \partsitem \( \begin{mat}[r]
                   1  &-1 &0  \\
                   0  &1  &-1 \\
                   0  &0  &1
                 \end{mat} \)
        \partsitem \( \begin{mat}[r]
                   1  &-1 &1/2  \\
                   1  &1  &-1/2 \\
                   0  &2  &0
                 \end{mat} \)
      \end{exparts*}  
      Sebagai contoh, untuk kolom pertama dari matriks pertama,
      $1=0\cdot x^2+1\cdot 1+0\cdot x$.
     \end{answer}
  \item 
    Untuk basis-basis dalam \nearbyexercise{exer:ChngeBasMatsRTwo},
    temukan matriks perubahan basis dalam arah sebaliknya, dari $D$ ke $B$.
    \begin{answer}
       Salah satu caranya adalah menemukan
       $\rep{\vec{\delta}_1}{B}$ dan $\rep{\vec{\delta}_2}{B}$,
       lalu menyusunnya sebagai kolom-kolom matriks perubahan basis yang
       diinginkan.
       Cara lain adalah mencari invers matriks-matriks yang menjawab
       \nearbyexercise{exer:ChngeBasMatsRTwo}. 
       \begin{exparts*}
        \partsitem
          $\begin{mat}[r]
            0  &1  \\
            1  &0
          \end{mat}$
        \partsitem \( \begin{mat}[r]
            1  &1  \\
            2  &4
          \end{mat} \)
        \partsitem \( \begin{mat}[r]
            2  &-1/2  \\
            -1 &1/2
          \end{mat} \)
        \partsitem \( \begin{mat}[r]
            2  &1  \\
            1  &1
          \end{mat} \)
      \end{exparts*} 
    \end{answer}
  \recommended \item 
    Tentukan apakah setiap matriks berikut mengubah basis pada \( \Re^2 \).
    Menjadi basis apakah \( \stdbasis_2 \) diubah?
    \begin{exparts*}
      \partsitem \( \begin{mat}[r]
                 5  &0  \\
                 0  &4
               \end{mat}  \)
      \partsitem \( \begin{mat}[r]
                 2  &1  \\
                 3  &1
               \end{mat}  \)
      \partsitem \( \begin{mat}[r]
                -1  &4  \\
                 2  &-8
               \end{mat}  \)
      \partsitem \( \begin{mat}[r]
                 1  &-1 \\
                 1  &1
               \end{mat}  \)
    \end{exparts*}
    \begin{answer}
       Suatu matriks mengubah basis jika dan hanya jika matriks tersebut
       nonsingular.
       \begin{exparts}
        \partsitem Matriks ini nonsingular sehingga mengubah basis.
          Menentukan menjadi basis apakah \( \stdbasis_2 \) diubah berarti
          mencari $D$ sedemikian sehingga
          \begin{equation*}
            \rep{\identity}{\stdbasis_2,D}=
            \begin{mat}[r]
                 5  &0  \\
                 0  &4
            \end{mat}
          \end{equation*}
          dan berdasarkan definisi representasi matriks bagi suatu pemetaan
          linear, kita memperoleh
          \begin{equation*}
            \rep{\identity(\vec{e}_1)}{D}=\rep{\vec{e}_1}{D}=\colvec[r]{5 \\ 0}
            \qquad
            \rep{\identity(\vec{e}_2)}{D}=\rep{\vec{e}_2}{D}=\colvec[r]{0 \\ 4}
          \end{equation*}
          Dengan
          \begin{equation*}
            D=\sequence{\colvec{x_1 \\ y_1},\colvec{x_2 \\ y_2}}
          \end{equation*}
          kita dapat menyelesaikan sistem
          \begin{equation*}
            \colvec[r]{1 \\ 0}
            =5\colvec{x_1 \\ y_1}+0\colvec{x_2 \\ y_2}
            \qquad
            \colvec[r]{0 \\ 1}
            =0\colvec{x_1 \\ y_1}+4\colvec{x_2 \\ y_2}
          \end{equation*}
          atau langsung mengenali jawabannya
          (dengan mengingat pembuktian \nearbylemma{le:NonSingIsChBasis}).
          \begin{equation*}
            D=\sequence{\colvec[r]{1/5 \\ 0},
                        \colvec[r]{0 \\ 1/4}}
          \end{equation*}
        \partsitem Ya, matriks ini nonsingular sehingga mengubah basis.
          Untuk menghitung $D$, kita melanjutkan seperti di atas dengan
          \begin{equation*}
            D=\sequence{\colvec{x_1 \\ y_1},
                        \colvec{x_2 \\ y_2}}
          \end{equation*}
          untuk menyelesaikan
          \begin{equation*}
            \colvec[r]{1 \\ 0}
            =2\colvec{x_1 \\ y_1}+3\colvec{x_2 \\ y_2}
            \quad\text{dan}\quad
            \colvec[r]{0 \\ 1}
            =1\colvec{x_1 \\ y_1}+1\colvec{x_2 \\ y_2}
          \end{equation*}
          dan memperoleh
          \begin{equation*}
            D=\sequence{\colvec[r]{-1 \\ 3},
                        \colvec[r]{1 \\ -2}}
          \end{equation*}
        \partsitem Tidak, matriks ini tidak mengubah basis karena singular.
        \partsitem Ya, matriks ini mengubah basis karena nonsingular.
          Perhitungan basis hasil perubahannya sama seperti di atas.
          \begin{equation*}
            D=\sequence{\colvec[r]{1/2 \\ -1/2},
                        \colvec[r]{1/2 \\ 1/2}}
          \end{equation*}
      \end{exparts}  
    \end{answer}
  \item Untuk setiap ruang, temukan matriks yang mengubah representasi vektor
    terhadap $B$ menjadi representasi terhadap~$D$.
    \begin{exparts}
      \partsitem $V=\Re^3$, $B=\stdbasis_3$, 
          $D=\sequence{\colvec{1 \\ 2 \\ 3},
                   \colvec{1 \\ 1 \\ 1},
                   \colvec{0 \\ 1 \\ -1}}$
      \partsitem $V=\Re^3$,  
        $B=\sequence{\colvec{1 \\ 2 \\ 3},
                     \colvec{1 \\ 1 \\ 1},
                     \colvec{0 \\ 1 \\ -1}}$,
        $D=\stdbasis_3$
      \partsitem $V=\polyspace_2$,
        $B=\sequence{x^2, x^2+x, x^2+x+1}$,
        $D=\sequence{2, -x, x^2}$
    \end{exparts}
    \begin{answer}
      \begin{exparts}
        \partsitem 
          Mulailah dengan menghitung pengaruh fungsi identitas pada setiap
          unsur basis awal~$B$. Jelas hasilnya adalah
          \begin{equation*}
            \colvec{1 \\ 0 \\ 0}\mapsunder{\identity}\colvec{1 \\ 0 \\ 0}
            \quad
            \colvec{0 \\ 1 \\ 0}\mapsunder{\identity}\colvec{0 \\ 1 \\ 0}
            \quad
            \colvec{0 \\ 0 \\ 1}\mapsunder{\identity}\colvec{0 \\ 0 \\ 1}
          \end{equation*}
          Sekarang representasikan ketiga keluaran itu terhadap basis akhir.
          \begin{equation*}
            \rep{\colvec{1 \\ 0 \\ 0}}{D}=\colvec{-2/3 \\ 5/3 \\ -1/3}
            \quad
            \rep{\colvec{0 \\ 1 \\ 0}}{D}=\colvec{1/3 \\ -1/3 \\ 2/3}
            \quad
            \rep{\colvec{0 \\ 0 \\ 1}}{D}=\colvec{1/3 \\ -1/3 \\ -1/3}
          \end{equation*}
          Susun ketiganya sebagai kolom-kolom matriks.
          \begin{equation*}
            \rep{\identity}{B,D}=
            \begin{mat}
              -2/3 &1/3  &1/3  \\
               5/3 &-1/3 &-1/3 \\
              -1/3 &2/3  &-1/3
            \end{mat}
          \end{equation*}
        \partsitem
          Salah satu cara menemukannya adalah mengambil invers matriks
          sebelumnya, karena matriks tersebut mengubah basis dalam arah
          sebaliknya. Sebagai alternatif, kita dapat menghitung ketiga
          representasi berikut:
          \begin{equation*}
            \rep{\colvec{1 \\ 2 \\ 3}}{\stdbasis_3}
              =\colvec{1 \\ 2 \\ 3}
            \quad
            \rep{\colvec{1 \\ 1 \\ 1}}{\stdbasis_3}
              =\colvec{1 \\ 1 \\ 1}
            \quad
            \rep{\colvec{0 \\ 1 \\ -1}}{\stdbasis_3}
               =\colvec{0 \\ 1 \\ -1}
          \end{equation*}
          lalu menyusunnya dalam sebuah matriks.
          \begin{equation*}
            \rep{\identity}{B,D}=
            \begin{mat}
              1 &1 &0 \\
              2 &1 &1 \\
              3 &1 &-1
            \end{mat}
          \end{equation*}
        \partsitem
          Merepresentasikan $\identity(x^2)$, $\identity(x^2+x)$,
          dan~$\identity(x^2+x+1)$ terhadap basis akhir menghasilkan
          \begin{equation*}
            \rep{x^2}{D}=\colvec{0 \\ 0 \\ 1}
            \quad
            \rep{x^2+x}{D}=\colvec{0 \\ -1 \\ 1}
            \quad
            \rep{x^2+x+1}{D}=\colvec{1/2 \\ -1 \\ 1}
          \end{equation*}
          Susun ketiganya menjadi satu matriks.
          \begin{equation*}
            \rep{\identity}{B,D}=
            \begin{mat}
              0 &0 &1/2 \\
              0 &-1 &-1 \\
              1 &1  &1
            \end{mat}
          \end{equation*}
      \end{exparts}
    \end{answer}
  \item 
    Temukan basis-basis sedemikian sehingga matriks berikut merepresentasikan
    pemetaan identitas terhadap basis-basis tersebut.
    \begin{equation*}
      \begin{mat}[r]
        3  &1  &4  \\
        2  &-1 &1  \\
        0  &0  &4
      \end{mat}
    \end{equation*}
    \begin{answer}
      Pertanyaan ini memiliki banyak penyelesaian berbeda.
      Salah satu caranya adalah memilih sebarang basis $B$ bagi sebarang ruang,
      lalu menghitung $D$ yang sesuai (tentu saja bagi ruang yang sama).
      Cara lain yang lebih mudah adalah menetapkan kodomain sebagai $\Re^3$
      dan basis kodomain sebagai $\stdbasis_3$.
      Dengan cara ini (ingat bahwa representasi sebarang vektor terhadap basis
      standar adalah vektor itu sendiri), kita memperoleh
      \begin{equation*}
        B=\sequence{\colvec[r]{3 \\ 2 \\ 0},
                    \colvec[r]{1 \\ -1 \\ 0},
                    \colvec[r]{4 \\ 1 \\ 4}  }
        \qquad
        D=\stdbasis_3
      \end{equation*}  
    \end{answer}
  \item 
    Tinjau ruang vektor fungsi-fungsi bernilai riil dengan basis
    \( \sequence{\sin(x),\cos(x)} \).
    Tunjukkan bahwa \( \sequence{2\sin(x)+\cos(x),3\cos(x)} \)
    juga merupakan basis bagi ruang ini.
    Temukan matriks perubahan basis dalam setiap arah.
    \begin{answer}
      Pemeriksaan bahwa
      \( B=\sequence{2\sin(x)+\cos(x),3\cos(x)} \) merupakan basis adalah
      rutin. Namai basis alaminya $D$.
      Untuk menghitung matriks perubahan basis $\rep{\identity}{B,D}$, kita
      harus menemukan $\rep{2\sin(x)+\cos(x)}{D}$ dan
      $\rep{3\cos(x)}{D}$; artinya, kita memerlukan $x_1,y_1,x_2,y_2$
      sedemikian sehingga persamaan-persamaan berikut berlaku.
      \begin{align*}
        x_1\cdot \sin(x)+y_1\cdot\cos(x) &= 2\sin(x)+\cos(x) \\
        x_2\cdot \sin(x)+y_2\cdot\cos(x) &= 3\cos(x) 
      \end{align*}
      Jelas jawabannya adalah
      \begin{equation*}
        \rep{\identity}{B,D}
        =\begin{mat}[r]
          2  &0  \\
          1  &3
        \end{mat}
      \end{equation*}
      Untuk matriks perubahan basis dalam arah sebaliknya, kita dapat mencari
      $\rep{\sin(x)}{B}$ dan $\rep{\cos(x)}{B}$ dengan menyelesaikan
      \begin{align*}
        w_1\cdot (2\sin(x)+\cos(x))+z_1\cdot(3\cos(x)) &= \sin(x) \\
        w_2\cdot (2\sin(x)+\cos(x))+z_2\cdot(3\cos(x)) &= \cos(x) 
      \end{align*}
      Cara yang lebih mudah adalah mencari invers matriks yang ditemukan di
      atas.
      \begin{equation*}
        \rep{\identity}{D,B}
        =\begin{mat}[r]
          2  &0  \\
          1  &3  
        \end{mat}^{-1}
        =\frac{1}{6}\cdot\begin{mat}[r]
          3  &0  \\
          -1  &2  
        \end{mat}
        =\begin{mat}[r]
           1/2  &0  \\
          -1/6  &1/3
         \end{mat}
      \end{equation*}
    \end{answer}
  \item 
    Basis standar \( \Re^2 \) dipetakan oleh matriks berikut ke basis yang
    mana?
    \begin{equation*}
        \begin{mat}
           \cos(2\theta)   &\sin(2\theta)   \\
           \sin(2\theta)   &-\cos(2\theta)
        \end{mat}
    \end{equation*}
    Bagaimana dengan basis-basis lainnya?
    \textit{Petunjuk.}
    Tinjau inversnya.
    \begin{answer} 
      Kita mulai dengan mengambil invers matriks tersebut, yaitu dengan
      menentukan pemetaan yang merupakan invers dari pemetaan yang dikaji.
      \begin{align*}
        \rep{\identity}{D,\stdbasis_2}
        =\rep{\identity}{\stdbasis_2,D}^{-1}
        &=\frac{1}{-\cos^2(2\theta)-\sin^2(2\theta)}\cdot\begin{mat}
           -\cos(2\theta)   &-\sin(2\theta)  \\
           -\sin(2\theta)   &\cos(2\theta)
        \end{mat}                                         \\
        &=\begin{mat}
           \cos(2\theta)   &\sin(2\theta)  \\
           \sin(2\theta)   &-\cos(2\theta)
        \end{mat}
      \end{align*}
      Representasi ini lebih mudah ditangani daripada representasi dalam arah
      sebaliknya karena matriks ini tersusun dari kedua vektor kolom berikut:
      \begin{equation*}
        \rep{\vec{\delta}_1}{\stdbasis_2}
           =\colvec{\cos(2\theta) \\ \sin(2\theta)}
        \qquad 
        \rep{\vec{\delta}_2}{\stdbasis_2}
           =\colvec{\sin(2\theta) \\ -\cos(2\theta)}
      \end{equation*}
      dan representasi terhadap $\stdbasis_2$ dapat langsung dibaca.
      \begin{equation*}
        \vec{\delta}_1=\colvec{\cos(2\theta) \\ \sin(2\theta)}
        \qquad 
        \vec{\delta}_2=\colvec{\sin(2\theta) \\ -\cos(2\theta)}
      \end{equation*}
      Gambar berikut memperlihatkan tindakan pemetaan yang mengubah $D$ menjadi
      $\stdbasis_2$ (sekali lagi, ini merupakan invers pemetaan yang menjadi
      jawaban pertanyaan tersebut).
      Garisnya membentuk sudut $\theta$ terhadap sumbu~$x$.
      \begin{center}  \small
        \includegraphics{map/mp/ch3.81}
      \end{center}
      Pemetaan ini mencerminkan vektor-vektor terhadap garis tersebut.
      Karena pencerminan merupakan invers dirinya sendiri, jawabannya
      adalah:~pemetaan semula mencerminkan terhadap garis yang melalui titik
      asal dan memiliki sudut elevasi $\theta$.
      (Tentu saja, pemetaan ini melakukan hal yang sama terhadap sebarang
      basis.)
    \end{answer}
  \recommended \item
    Apa matriks perubahan basis terhadap \( B,B \)?
    \begin{answer}
       Matriks identitas dengan ukuran yang sesuai.
     \end{answer}
  \item 
    Buktikan bahwa suatu matriks mengubah basis jika dan hanya jika matriks
    tersebut dapat diinvers.
    \begin{answer}
       Kedua sifat tersebut berlaku jika dan hanya jika matriksnya nonsingular.
    \end{answer}
  \item 
    Selesaikan pembuktian \nearbylemma{le:NonSingIsChBasis}.
    \begin{answer}
      Yang tersisa adalah menunjukkan bahwa perkalian kiri dengan matriks
      reduksi merepresentasikan perubahan dari suatu basis lain ke
      \( B=\basis{\beta}{n} \).

      Penerapan matriks perkalian baris \( M_i(k) \) mengubah representasi
      terhadap basis
      \( \sequence{\vec{\beta}_1,\dots,k\vec{\beta}_i,\dots,\vec{\beta}_n} \)
      menjadi representasi terhadap \( B \), seperti berikut.
      \begin{equation*}
         \vec{v}=c_1\cdot\vec{\beta}_1+\dots+c_i\cdot(k\vec{\beta}_i)
                   +\dots+c_n\cdot\vec{\beta}_n  
         \;\mapsto\;                                                       
         c_1\cdot\vec{\beta}_1+\dots
             +(kc_i)\cdot\vec{\beta}_i+\dots+c_n\cdot\vec{\beta}_n=\vec{v}
      \end{equation*}
      Penerapan matriks pertukaran baris \( P_{i,j} \) mengubah representasi
      terhadap basis
      \( \sequence{\vec{\beta}_1,\dots,\vec{\beta}_j,\dots,
        \vec{\beta}_i,\dots,\vec{\beta}_n} \)
      menjadi representasi terhadap
            \( \sequence{\vec{\beta}_1,\dots,\vec{\beta}_i,\dots,
        \vec{\beta}_j,\dots,\vec{\beta}_n} \).
      Terakhir, penerapan matriks kombinasi baris \( C_{i,j}(k) \) mengubah
      representasi terhadap
      \( \sequence{\vec{\beta}_1,\dots,\vec{\beta}_i+k\vec{\beta}_j,\dots,
        \vec{\beta}_j,\dots,\vec{\beta}_n} \)
      menjadi representasi terhadap \( B \).
      \begin{multline*}
         \vec{v}= c_1\cdot\vec{\beta}_1+\dots
                     +c_i\cdot(\vec{\beta}_i+k\vec{\beta}_j)
                       +\dots+c_j\vec{\beta}_j+\dots+c_n\cdot\vec{\beta}_n  
         \\  \mapsto\;                                                       
         c_1\cdot\vec{\beta}_1+\dots+c_i\cdot\vec{\beta}_i
               +\dots+(kc_i+c_j)\cdot\vec{\beta}_j
               +\dots+c_n\cdot\vec{\beta}_n=\vec{v}     
      \end{multline*}
      (Seperti pada bagian pembuktian dalam isi subbagian ini, berbagai syarat
      pada operasi baris, misalnya bahwa skalar $k$ tak nol, memastikan bahwa
      semua barisan tersebut benar-benar merupakan basis.)
    \end{answer}
  \recommended \item 
    Misalkan \( H \) suatu matriks nonsingular \( \nbyn{n} \).
    Basis manakah dari \( \Re^n \) yang diubah oleh \( H \) menjadi basis
    standar?
    \begin{answer}
       Dengan mengambil $H$ sebagai matriks perubahan basis
       $H=\rep{\identity}{B,\stdbasis_n}$, 
       kolom-kolomnya adalah
       \begin{equation*}
         \colvec{h_{1,i} \\ \vdots \\ h_{n,i}}
         =\rep{\identity(\vec{\beta}_i)}{\stdbasis_n}
         =\rep{\vec{\beta}_i}{\stdbasis_n}
       \end{equation*}
       dan karena representasi terhadap basis standar dapat langsung dibaca,
       kita memperoleh
       \begin{equation*}
         \colvec{h_{1,i} \\ \vdots \\ h_{n,i}}
         =\vec{\beta}_i
       \end{equation*}
       Jadi, basis tersebut terdiri atas kolom-kolom \( H \).
    \end{answer}
  \recommended \item \label{exer:AnyNonZeroRepChgTOAnyOther}
    \begin{exparts}
      \partsitem Dalam \( \polyspace_3 \) dengan basis
        \( B=\sequence{1+x,1-x,x^2+x^3,x^2-x^3} \), kita memiliki
        representasi berikut.
        \begin{equation*}
          \rep{1-x+3x^2-x^3}{B}=
            \colvec[r]{0 \\ 1 \\ 1 \\ 2}_B
        \end{equation*}
        Temukan basis \( D \) yang memberikan representasi berbeda berikut
        bagi polinomial yang sama.
        \begin{equation*}
          \rep{1-x+3x^2-x^3}{D}=
            \colvec[r]{1 \\ 0 \\ 2 \\ 0}_D
        \end{equation*}
      \partsitem Nyatakan dan buktikan bahwa kita dapat mengubah sebarang
        representasi vektor tak nol menjadi sebarang representasi lainnya.
    \end{exparts}
    \noindent\textit{Petunjuk.}
    Pembuktian \nearbylemma{le:NonSingIsChBasis} bersifat konstruktif\Dash
    pembuktian itu tidak sekadar menyatakan bahwa basis berubah, tetapi juga
    memperlihatkan bagaimana basis tersebut berubah.
    \begin{answer}
      \begin{exparts}
        \partsitem Kita dapat mengubah representasi vektor awal menjadi
          representasi akhir melalui serangkaian operasi baris.
          Pembuktian tersebut menunjukkan bagaimana basis-basisnya berubah.
          Kita mulai dengan mempertukarkan baris pertama dan kedua dari
          representasi terhadap $B$ untuk memperoleh representasi terhadap
          basis baru $B_1$.
          \begin{equation*}
            \rep{1-x+3x^2-x^3}{B_1}=
              \colvec[r]{1 \\ 0 \\ 1 \\ 2}_{B_1}
            \qquad
            B_1=\sequence{1-x,1+x,x^2+x^3,x^2-x^3}
          \end{equation*}
          Selanjutnya, kita menambahkan \( -2 \) kali baris ketiga
          representasi vektor ke baris keempat.
          \begin{equation*}
            \rep{1-x+3x^2-x^3}{B_2}=
              \colvec[r]{1 \\ 0 \\ 1 \\ 0}_{B_2}
            \qquad
            B_2=\sequence{1-x,1+x,3x^2-x^3,x^2-x^3}
          \end{equation*}
          (Unsur ketiga \( B_2 \) adalah unsur ketiga \( B_1 \) dikurangi
          \( -2 \) kali unsur keempat $B_1$.)
          Sekarang kita dapat menyelesaikannya dengan menggandakan baris ketiga.
          \begin{equation*}
            \rep{1-x+3x^2-x^3}{D}=
              \colvec[r]{1 \\ 0 \\ 2 \\ 0}_{D}
            \qquad
            D=\sequence{1-x,1+x,(3x^2-x^3)/2,x^2-x^3}
          \end{equation*}
        \partsitem 
          Berikut tiga cara berbeda untuk menyatakan hasil semacam itu.
          Pernyataan pertama adalah:~jika $V$ merupakan ruang vektor dengan
          basis $B$ dan $\vec{v}\in V$ tak nol, maka untuk setiap vektor kolom
          tak nol $\vec{z}$ (yang banyak komponennya sama dengan dimensi $V$),
          terdapat matriks perubahan basis $M$ sedemikian sehingga
          $M\cdot \rep{\vec{v}}{B}=\vec{z}$. 
          Pernyataan kedua yang mungkin adalah:~untuk sebarang ruang vektor
          $V$ berdimensi $n$ dan sebarang vektor tak nol \( \vec{v}\in V \),
          jika \( \vec{z}_1, \vec{z}_2\in\Re^n \) tak nol, maka terdapat basis
          \( B,D\subset V \) sedemikian sehingga
          \( \rep{\vec{v}}{B}=\vec{z}_1 \) dan
          \( \rep{\vec{v}}{D}=\vec{z}_2 \).
          Pernyataan ketiga adalah:~untuk sebarang anggota tak nol $\vec{v}$
          dari sebarang ruang vektor berdimensi~$n$ dan sebarang vektor kolom
          tak nol dengan $n$ komponen, terdapat basis sedemikian sehingga
          $\vec{v}$ direpresentasikan terhadap basis tersebut oleh vektor
          kolom itu.

          Pernyataan pertama dan kedua mudah mengikuti pernyataan ketiga.
          Pernyataan pertama mengikuti karena pernyataan ketiga memberikan
          basis $D$ sedemikian sehingga $\rep{\vec{v}}{D}=\vec{z}$, lalu
          $\rep{\identity}{B,D}$ merupakan~$M$ yang diinginkan.
          Pernyataan kedua mengikuti pernyataan ketiga karena hanya merupakan
          penerapannya sebanyak dua kali.

          Salah satu cara membuktikan pernyataan ketiga sama seperti jawaban
          bagian pertama pertanyaan ini. Berikut sketsanya.
          Representasikan $\vec{v}$ terhadap sebarang basis $B$ dengan vektor
          kolom $\vec{z}_1$.
          Vektor kolom ini harus memiliki suatu komponen tak nol karena
          $\vec{v}$ merupakan vektor tak nol.
          Gunakan komponen tersebut dalam serangkaian operasi baris untuk
          mengubah $\vec{z}_1$ menjadi $\vec{z}$.
          (Sketsa ini dapat kita lengkapi menjadi argumen induksi pada dimensi
          $V$.)
       \end{exparts}  
     \end{answer}
  \item 
    Misalkan \( V,W \) ruang-ruang vektor, \( B,\hat{B} \) basis-basis bagi
    \( V \), dan \( D,\hat{D} \) basis-basis bagi \( W \).
    Jika \( \map{h}{V}{W} \) linear, temukan rumus yang menghubungkan
    \( \rep{h}{B,D} \) dengan \( \rep{h}{\hat{B},\hat{D}} \).
    \begin{answer}
      Inilah topik subbagian berikutnya.
    \end{answer}
  \recommended \item
    Tunjukkan bahwa kolom-kolom matriks perubahan basis \( \nbyn{n} \)
    membentuk basis bagi \( \Re^n \).
    Apakah semua basis muncul dengan cara itu:~dapatkah vektor-vektor dari
    sebarang basis $\Re^n$ menjadi kolom-kolom suatu matriks perubahan basis?
    \begin{answer}
      Matriks perubahan basis bersifat nonsingular sehingga ranknya sama dengan
      banyak kolomnya. Oleh karena itu, himpunan kolomnya merupakan subhimpunan
      bebas linear berukuran $n$ dalam $\Re^n$, sehingga merupakan basis.
      Jawaban bagi bagian kedua juga `ya'; semua implikasi dalam kalimat
      sebelumnya berlaku dalam arah sebaliknya (artinya, semua bagian
      `jika \ldots maka~\ldots' pada kalimat sebelumnya berubah menjadi
      `jika dan hanya jika').
    \end{answer}
  \recommended \item 
    Temukan suatu matriks yang memiliki pengaruh berikut.
    \begin{equation*}
      \colvec[r]{1 \\ 3}
      \;\mapsto\;
      \colvec[r]{4 \\ -1}
    \end{equation*}
    Artinya, temukan $M$ yang melalui perkalian kiri mengubah vektor awal
    menjadi vektor akhir.
    Adakah matriks yang memiliki kedua pengaruh berikut?
    \begin{exparts*}
      \partsitem
        $
          \colvec[r]{1 \\ 3}\mapsto\colvec[r]{1 \\ 1}
          \quad
          \colvec[r]{2 \\ -1}\mapsto\colvec[r]{-1 \\ -1}
        $
      \partsitem $
          \colvec[r]{1 \\ 3}\mapsto\colvec[r]{1 \\ 1}
          \quad
          \colvec[r]{2 \\ 6}\mapsto\colvec[r]{-1 \\ -1}
        $
   \end{exparts*}
   Berikan syarat perlu dan cukup agar terdapat matriks sedemikian sehingga
   $\vec{v}_1\mapsto\vec{w}_1$ dan $\vec{v}_2\mapsto\vec{w}_2$.
    \begin{answer}
      Untuk bagian pertama pertanyaan, terdapat tak hingga banyak matriks
      semacam itu. Salah satunya merepresentasikan terhadap \( \stdbasis_2 \)
      transformasi pada \( \Re^2 \) dengan tindakan berikut.
      \begin{equation*}
        \colvec[r]{1 \\ 0}\mapsto\colvec[r]{4 \\ 0}
        \qquad
        \colvec[r]{0 \\ 1}\mapsto\colvec[r]{0 \\ -1/3}
      \end{equation*}  
     Persoalan menetapkan dua pasangan masukan/keluaran berbeda sedikit lebih
     rumit. Fakta bahwa matriks bertindak secara linear meniadakan beberapa
     kemungkinan.
     \begin{exparts}
       \partsitem Ya, terdapat matriks semacam itu. Syarat-syarat berikut
         \begin{equation*}
           \begin{mat}
             a  &b  \\
             c  &d
           \end{mat}
           \colvec[r]{1 \\ 3}
           =
           \colvec[r]{1 \\ 1}
           \qquad
           \begin{mat}
             a  &b  \\
             c  &d
           \end{mat}
           \colvec[r]{2 \\ -1}
           =
           \colvec[r]{-1 \\ -1}
         \end{equation*}
         dapat diselesaikan
         \begin{equation*}
           \begin{linsys}{4}
             a  &+  &3b  &   &   &   &   &=  &1  \\
                &   &    &   &c  &+  &3d &=  &1  \\
            2a  &-  &b   &   &   &   &   &=  &-1 \\
                &   &    &   &2c &-  &d  &=  &-1 
           \end{linsys}
         \end{equation*}
         untuk menghasilkan matriks berikut.
         \begin{equation*}
           \begin{mat}[r]
             -2/7  &3/7 \\
             -2/7  &3/7 
           \end{mat}
         \end{equation*}
       \partsitem Tidak, karena
         \begin{equation*}
           2\cdot\colvec[r]{1 \\ 3}=\colvec[r]{2 \\ 6}
           \quad\text{tetapi}\quad
           2\cdot\colvec[r]{1 \\ 1}\neq\colvec[r]{-1 \\ -1}
         \end{equation*}
         sehingga tidak ada tindakan linear yang dapat menghasilkan pengaruh
         ini.
       \partsitem Salah satu syarat cukup adalah bahwa
         \( \set{\vec{v}_1,\vec{v}_2} \) bebas linear, tetapi syarat ini tidak
         perlu. Syarat perlu dan cukupnya adalah bahwa setiap kebergantungan
         linear di antara vektor-vektor awal juga muncul di antara
         vektor-vektor akhir. Artinya,
         \begin{equation*}
           c_1\vec{v}_1+c_2\vec{v}_2=\zero
           \quad\text{menyiratkan}\quad
           c_1\vec{w}_1+c_2\vec{w}_2=\zero.
         \end{equation*}
         Pembuktian syarat ini bersifat rutin.
     \end{exparts} 
   \end{answer}
\end{exercises}

















\subsection{Mengubah Representasi Pemetaan}
\index{ekuivalensi matriks|(}
Subbagian pertama menunjukkan cara mengubah representasi suatu vektor terhadap
satu basis menjadi representasi vektor yang sama terhadap basis lain.
Selanjutnya, kita mengubah representasi suatu pemetaan terhadap satu pasangan
basis menjadi representasi terhadap pasangan lain\Dash kita mengubah
$\rep{h}{B,D}$ menjadi $\rep{h}{\hat{B},\hat{D}}$.
Berikut diagram panahnya.\index{diagram panah}
%<*ChangeRepresentationOfMapArrowDiagram>
\begin{equation*}
  \begin{CD}
    V_{\wrt{B}}                   @>h>H>        W_{\wrt{D}}       \\
    @V{\text{\scriptsize$\identity$}} VV                @V{\text{\scriptsize$\identity$}} VV \\
    V_{\wrt{\hat{B}}}             @>h>\hat{H}>  W_{\wrt{\hat{D}}}
  \end{CD}
\end{equation*}
Untuk bergerak dari kiri bawah ke kanan bawah, kita dapat langsung bergerak
mendatar, atau naik ke $V_B$, bergerak mendatar ke $W_D$, lalu turun.
Jadi, kita dapat menghitung $\hat{H}=\rep{h}{\hat{B},\hat{D}}$ secara langsung
dengan menggunakan $\hat{B}$ dan $\hat{D}$, atau terlebih dahulu mengubah basis
dengan $\rep{\identity}{\hat{B},B}$, kemudian mengalikan dengan
\( H=\rep{h}{B,D} \), lalu mengubah basis dengan
$\rep{\identity}{D,\hat{D}}$.
%</ChangeRepresentationOfMapArrowDiagram>

\begin{theorem}
%<*ChangeRepresentationOfMapTheorem>
Untuk mengubah matriks~$H$ yang merepresentasikan pemetaan~$h$ terhadap $B,D$
menjadi matriks~$\hat{H}$ yang merepresentasikannya terhadap
$\hat{B},\hat{D}$, gunakan rumus berikut.
\begin{equation*}
   \hat{H}=
   \rep{\identity}{D,\hat{D}}\cdot H\cdot \rep{\identity}{\hat{B},B}
   \tag*{\text{($*$})}
\end{equation*}
%</ChangeRepresentationOfMapTheorem>
\end{theorem}

\begin{proof}
%<*ChangeRepresentationOfMapTheoremPf>
Hal ini langsung tampak dari diagram tersebut.
%</ChangeRepresentationOfMapTheoremPf>
\end{proof}

% (To compare the equation with the sentence before it
% remember to read its right hand side
% from right to left, because we read
% function composition from right to left and matrix multiplication 
% represents composition).

\begin{example}
Matriks
\begin{equation*}
  T=\begin{mat}[r]
     \cos(\pi/6)  &-\sin(\pi/6)  \\
     \sin(\pi/6)  &\cos(\pi/6)
  \end{mat}
  =
  \begin{mat}[r]
     \sqrt{3}/2  &-1/2  \\
     1/2         &\sqrt{3}/2
  \end{mat}
\end{equation*}
merepresentasikan, terhadap \( \stdbasis_2,\stdbasis_2 \), transformasi
\( \map{t}{\Re^2}{\Re^2} \) yang memutar vektor-vektor berlawanan arah jarum
jam sebesar sudut \( \pi/6 \) radian.
\begin{center}
  \includegraphics{map/mp/ch3.22}
\end{center}
Kita dapat mengubah $T$ menjadi representasi terhadap basis-basis berikut:
\begin{equation*}
  \hat{B}=\sequence{
              \colvec[r]{1 \\ 1},
              \colvec[r]{0 \\ 2} }
  \qquad
  \hat{D}=\sequence{
              \colvec[r]{-1 \\ 0},
              \colvec[r]{2 \\ 3} }
\end{equation*}
dengan menggunakan diagram panah di atas.
\begin{equation*}
  \begin{CD}
    \Re^2_{\wrt{\stdbasis_2}}              @>t>T>        \Re^2_{\wrt{\stdbasis_2}}     \\
    @V\text{\scriptsize$\identity$} VV                @V\text{\scriptsize$\identity$} VV \\
    \Re^2_{\wrt{\hat{B}}}          @>t>\hat{T}>  \Re^2_{\wrt{\hat{D}}}
  \end{CD}
\end{equation*}
Gambar tersebut memperlihatkan bahwa kita dapat menghitung~$\hat{T}$ secara
langsung dengan bergerak sepanjang sisi bawah persegi, atau seperti dalam
rumus~($*$), dengan naik pada sisi kiri, bergerak sepanjang sisi atas, lalu
turun pada sisi kanan, dengan
$\hat{T}=
   \rep{\identity}{\stdbasis_2,\hat{D}}\cdot T\cdot \rep{\identity}{\hat{B},\stdbasis_2}$.
(Perhatikan sekali lagi bahwa perkalian matriks dibaca dari kanan ke kiri,
karena ketiga fungsi dikomposisikan dan komposisi fungsi dibaca dari kanan ke
kiri.)

Temukan matriks untuk sisi kiri diagram,
$\rep{\identity}{\hat{B},\stdbasis_2}$, dengan cara biasa:~temukan pengaruh
pemetaan identitas pada basis awal~$\hat{B}$\Dash yakni sama sekali tidak ada
perubahan\Dash lalu representasikan unsur-unsur basis tersebut terhadap basis
akhir~$\stdbasis_2$.
\begin{equation*}
  \rep{\identity}{\hat{B},\stdbasis_2}
  =
  \begin{mat}
    1 &0 \\
    1 &2
  \end{mat}
\end{equation*}
Perhitungan ini mudah jika basis akhirnya merupakan basis standar.

Terdapat dua cara untuk menghitung matriks bagi gerak turun pada sisi kanan
persegi, $\rep{\identity}{\stdbasis_2,\hat{D}}$.
Kita dapat menghitungnya secara langsung seperti pada matriks perubahan basis
lainnya. Sebagai alternatif, kita dapat menghitungnya sebagai invers matriks
bagi gerak naik, $\rep{\identity}{\hat{D},\stdbasis_2}$.
Matriks tersebut mudah ditemukan dan kita memiliki rumus invers $\nbyn{2}$,
sehingga itulah yang digunakan dalam persamaan berikut.
\begin{align*}
   \rep{t}{\hat{B},\hat{D}}
  &=\begin{mat}[r]
     -1     &2   \\
     0      &3
  \end{mat}^{-1}
  \begin{mat}[r]
     \sqrt{3}/2  &-1/2  \\
     1/2         &\sqrt{3}/2
  \end{mat}
  \begin{mat}[r]
     1      &0   \\
     1      &2
  \end{mat}                              \\                            
  &=\begin{mat}[r]
     (5-\sqrt{3})/6   &(3+2\sqrt{3})/3 \\
     (1+\sqrt{3})/6  &\sqrt{3}/3
  \end{mat}
\end{align*}

Matriksnya lebih rumit, tetapi pemetaan yang direpresentasikannya tetap sama.
Sebagai contoh, untuk menghasilkan kembali pengaruh $t$ pada gambar, mulailah
dengan $\hat{B}$,
\begin{equation*}
  \rep{\colvec[r]{1 \\ 3}}{\hat{B}}=\colvec[r]{1 \\ 1}_{\hat{B}}
\end{equation*}
terapkan $\hat{T}$,
\begin{equation*}
  \begin{mat}[r]
     (5-\sqrt{3})/6   &(3+2\sqrt{3})/3 \\
     (1+\sqrt{3})/6  &\sqrt{3}/3
  \end{mat}_{\hat{B},\hat{D}}
  \colvec[r]{1 \\ 1}_{\hat{B}}
  =
  \colvec[r]{(11+3\sqrt{3})/6 \\ (1+3\sqrt{3})/6}_{\hat{D}}
\end{equation*}
dan periksa hasilnya terhadap $\hat{D}$.
\begin{equation*}
  \frac{11+3\sqrt{3}}{6}\cdot\colvec[r]{-1 \\ 0}
  +\frac{1+3\sqrt{3}}{6}\cdot\colvec[r]{2 \\ 3}
  =\colvec[r]{(-3+\sqrt{3})/2 \\ (1+3\sqrt{3})/2}
\end{equation*}
\end{example}

\begin{example} \label{ex:DiagizedMat}
Mengubah basis dapat menyederhanakan matriks.
Pada \( \Re^3 \), pemetaan
\begin{equation*}
  \colvec{x \\ y \\ z}\mapsunder{t}\colvec{y+z \\ x+z \\ x+y}
\end{equation*}
direpresentasikan terhadap basis standar sebagai berikut.
\begin{equation*}
  \rep{t}{\stdbasis_3,\stdbasis_3}=
  \begin{mat}[r]
    0  &1  &1  \\
    1  &0  &1  \\
    1  &1  &0
  \end{mat}
\end{equation*}
Merepresentasikannya terhadap
\begin{equation*}
  B=\sequence{\colvec[r]{1 \\ -1 \\ 0},
                             \colvec[r]{1 \\ 1 \\ -2},
                             \colvec[r]{1 \\ 1 \\ 1}}
\end{equation*} 
menghasilkan matriks diagonal.
\begin{equation*} 
  \rep{t}{B,B}=
  \begin{mat}[r]
   -1  &0  &0  \\
    0  &-1 &0  \\
    0  &0  &2
  \end{mat}
\end{equation*}
\end{example}

Tentu saja, biasanya kita lebih menyukai representasi yang lebih mudah
dipahami. Kita mengatakan bahwa suatu pemetaan atau matriks telah
\definend{didiagonalkan}\index{matriks!didiagonalkan} jika kita menemukan
basis~$B$ sedemikian sehingga representasinya diagonal terhadap $B,B$,
yakni terhadap basis awal dan basis akhir yang sama.
Bab Lima menentukan pemetaan dan matriks mana yang dapat didiagonalkan.

Sisa subbagian ini mengembangkan kasus yang lebih mudah, yaitu menemukan dua
basis $B,D$ sedemikian sehingga suatu representasi menjadi sederhana.
Ingat bahwa subbagian sebelumnya menunjukkan bahwa suatu matriks merupakan
matriks perubahan basis jika dan hanya jika matriks tersebut nonsingular.
% That gives us another version of the above  arrow diagram
% and equation~($*$) from the start of this subsection.

\begin{definition}  \label{def:MatEquiv}
%<*df:MatEquiv>
Matriks-matriks berukuran sama \( H \) dan \( \hat{H} \) disebut
\definend{ekuivalen secara matriks\/}\index{ekuivalensi matriks!definisi}%
\index{matriks!ekuivalen}\index{relasi ekuivalensi!ekuivalensi matriks}
jika terdapat matriks-matriks nonsingular \( P \) dan \( Q \) sedemikian
sehingga
$\hat{H}=PHQ$.
%</df:MatEquiv>
\end{definition}

\begin{corollary} \label{le:MatEqIsSameMap}
%<*co:MatEqIsSameMap>
Matriks-matriks yang ekuivalen secara matriks merepresentasikan pemetaan yang
sama terhadap pasangan basis yang sesuai.
%</co:MatEqIsSameMap>
\end{corollary}

\begin{proof}
Hal ini langsung mengikuti persamaan~($*$) di atas.
\end{proof}

\nearbyexercise{exer:MatEqIsEqRel} memeriksa bahwa ekuivalensi matriks
merupakan relasi ekuivalensi.
Dengan demikian, relasi ini mempartisi
\index{partisi!kelas ekuivalensi matriks} himpunan matriks menjadi
kelas-kelas ekuivalensi matriks.
\begin{center}
  \raisebox{.5in}{\begin{tabular}{l}
                    Semua matriks:
                  \end{tabular}}
  \includegraphics{map/mp/ch3.23}
  \raisebox{.3in}{\begin{tabular}{l}
                     $H$ ekuivalen secara matriks \\ dengan $\hat{H}$
                  \end{tabular}}
\end{center}
Kita dapat memahami kelas-kelas tersebut dengan membandingkan ekuivalensi
matriks dengan ekuivalensi baris (ingat bahwa dua matriks ekuivalen baris jika
keduanya dapat direduksi satu sama lain melalui operasi baris).
Dalam $\hat{H}=PHQ$, matriks $P$ dan $Q$ nonsingular sehingga, menurut
Lema~IV.\ref{lem:ComputeInvMat}, masing-masing merupakan hasil kali
matriks-matriks reduksi elementer.
Perkalian kiri dengan matriks-matriks reduksi yang membentuk $P$ menjalankan
operasi baris. Perkalian kanan dengan matriks-matriks reduksi yang membentuk
$Q$ menjalankan operasi kolom.
Jadi, ekuivalensi matriks merupakan perumuman ekuivalensi baris\Dash dua
matriks ekuivalen baris jika salah satunya dapat diubah menjadi yang lain
melalui serangkaian langkah reduksi baris, sedangkan dua matriks ekuivalen
secara matriks jika salah satunya dapat diubah menjadi yang lain melalui
serangkaian langkah reduksi baris yang diikuti serangkaian langkah reduksi
kolom.

Akibatnya, jika dua matriks ekuivalen baris, keduanya juga ekuivalen secara
matriks karena kita dapat mengambil $Q$ sebagai matriks identitas.
Namun, kebalikannya tidak berlaku:~dua matriks dapat ekuivalen secara matriks
tanpa ekuivalen baris.

\begin{example}
Kedua matriks berikut ekuivalen secara matriks
\begin{equation*} 
  \begin{mat}[r]
    1  &0  \\
    0  &0
  \end{mat}
  \qquad
  \begin{mat}[r]
    1  &1  \\
    0  &0
  \end{mat}
\end{equation*}
karena matriks kedua tereduksi menjadi matriks pertama melalui operasi kolom
yang menambahkan $-1$ kali kolom pertama ke kolom kedua.
Keduanya tidak ekuivalen baris karena memiliki bentuk eselon tereduksi yang
berbeda (keduanya sudah berada dalam bentuk tereduksi).
\end{example}

Kita menutup bagian ini dengan memberikan suatu himpunan wakil bagi
kelas-kelas ekuivalensi matriks. % \appendrefs{class representatives}%
\index{wakil!kelas ekuivalensi matriks}

\begin{theorem}  \label{th:CanonFormForMatEquiv}
\index{ekuivalensi matriks!bentuk kanonik}
\index{bentuk kanonik!ekuivalensi matriks}
%<*th:CanonFormForMatEquiv>
Sebarang matriks \( \nbym{m}{n} \) berank \( k \) ekuivalen secara matriks
dengan matriks \( \nbym{m}{n} \) yang semua entrinya nol kecuali \( k \)
entri diagonal pertamanya yang bernilai satu.
\begin{equation*}
    \begin{mat}
      1  &0      &\ldots &0  &0  &\ldots  &0  \\
      0  &1      &\ldots &0  &0  &\ldots  &0  \\
         &\vdots                              \\
      0  &0      &\ldots &1  &0  &\ldots  &0  \\
      0  &0      &\ldots &0  &0  &\ldots  &0  \\
         &\vdots                              \\
      0  &0      &\ldots &0  &0  &\ldots  &0
    \end{mat}
\end{equation*}
%</th:CanonFormForMatEquiv>
\end{theorem}

%<*BlockPartialIdentityForm>
\noindent Ini merupakan bentuk
\definend{blok identitas parsial}\index{matriks!blok}.
\begin{equation*}
    \begin{pmat}{c|c}
      I  &Z  \\  \hline
      Z  &Z
    \end{pmat}
\end{equation*}
%</BlockPartialIdentityForm>

\begin{proof}
%<*pf:CanonFormForMatEquiv>
Reduksi matriks yang diberikan dengan metode Gauss--Jordan dan gabungkan semua
matriks reduksi baris untuk membentuk \( P \).
Kemudian gunakan entri-entri utama untuk melakukan reduksi kolom dan akhiri
dengan mempertukarkan kolom-kolom agar satu-satu utama berada pada diagonal.
Gabungkan matriks-matriks reduksi kolom menjadi \( Q \).
%</pf:CanonFormForMatEquiv>
\end{proof}

\begin{example}
Kita mengilustrasikan pembuktian dengan menemukan $P$ dan $Q$ bagi matriks
berikut.
\begin{equation*}
    \begin{mat}[r]
       1  &2  &1  &-1  \\
       0  &0  &1  &-1  \\
       2  &4  &2  &-2
    \end{mat}
\end{equation*}
Pertama, lakukan reduksi baris Gauss--Jordan.
\begin{equation*}
    \begin{mat}[r]
       1  &-1 &0    \\
       0  &1  &0    \\
       0  &0  &1
    \end{mat}
    \begin{mat}[r]
       1  &0  &0    \\
       0  &1  &0    \\
       -2 &0  &1
    \end{mat}
    \begin{mat}[r]
       1  &2  &1  &-1  \\
       0  &0  &1  &-1  \\
       2  &4  &2  &-2
    \end{mat}
  =
    \begin{mat}[r]
       1  &2  &0  &0   \\
       0  &0  &1  &-1  \\
       0  &0  &0  &0
    \end{mat}
\end{equation*}
Kemudian lakukan reduksi kolom, yang melibatkan perkalian kanan.
\begin{equation*}
    \begin{mat}[r]
       1  &2  &0  &0   \\
       0  &0  &1  &-1  \\
       0  &0  &0  &0
    \end{mat}
    \begin{mat}[r]
       1  &-2 &0  &0   \\
       0  &1  &0  &0   \\
       0  &0  &1  &0   \\
       0  &0  &0  &1
    \end{mat}
    \begin{mat}[r]
       1  &0  &0  &0   \\
       0  &1  &0  &0   \\
       0  &0  &1  &1   \\
       0  &0  &0  &1
    \end{mat}
  =
    \begin{mat}[r]
       1  &0  &0  &0   \\
       0  &0  &1  &0  \\
       0  &0  &0  &0
    \end{mat}
\end{equation*}
Akhiri dengan mempertukarkan kolom.
\begin{equation*}
    \begin{mat}[r]
       1  &0  &0  &0   \\
       0  &0  &1  &0  \\
       0  &0  &0  &0
    \end{mat}
    \begin{mat}[r]
       1  &0  &0  &0   \\
       0  &0  &1  &0   \\
       0  &1  &0  &0   \\
       0  &0  &0  &1
    \end{mat}
  =
    \begin{mat}[r]
       1  &0  &0  &0   \\
       0  &1  &0  &0  \\
       0  &0  &0  &0
    \end{mat}
\end{equation*}
Terakhir, gabungkan semua pengali kiri menjadi $P$ dan semua pengali kanan
menjadi $Q$ untuk memperoleh $PHQ$.
\begin{equation*}
    \begin{mat}[r]
       1  &-1 &0    \\
       0  &1  &0    \\
       -2 &0  &1
    \end{mat}
    \begin{mat}[r]
       1  &2  &1  &-1  \\
       0  &0  &1  &-1  \\
       2  &4  &2  &-2
    \end{mat}
    \begin{mat}[r]
       1  &0  &-2  &0   \\
       0  &0  &1  &0   \\
       0  &1  &0  &1   \\
       0  &0  &0  &1
    \end{mat}
    =
    \begin{mat}[r]
       1  &0  &0  &0   \\
       0  &1  &0  &0  \\
       0  &0  &0  &0
    \end{mat}
\end{equation*}
\end{example}

\begin{corollary}  \label{co:MatrixEquivalentIffSameRank}
%<*co:MatrixEquivalentIffSameRank>
Kelas-kelas ekuivalensi matriks dicirikan\index{mencirikan} oleh rank:~dua
matriks berukuran sama ekuivalen secara matriks jika dan hanya jika keduanya
memiliki rank yang sama.
%</co:MatrixEquivalentIffSameRank>
\end{corollary}

\begin{proof}
%<*pf:MatrixEquivalentIffSameRank>
Dua matriks berukuran sama dengan rank yang sama ekuivalen dengan matriks blok
identitas parsial yang sama.
%</pf:MatrixEquivalentIffSameRank>
\end{proof}

\begin{example}
Matriks-matriks $\nbyn{2}$ hanya memiliki tiga kemungkinan rank:~nol, satu,
atau~dua. Jadi, terdapat tiga kelas ekuivalensi matriks.
\index{partisi!kelas ekuivalensi matriks}
\begin{center} % make an exercise comparing these classes with the ones for row-equivalence?
  \raisebox{.5in}{\begin{tabular}{l}
                    Semua matriks $\nbyn{2}$:
                  \end{tabular}}
  \includegraphics{map/mp/ch3.24}
  \raisebox{.3in}{\begin{tabular}{l}
                    Tiga kelas \\ ekuivalensi
                  \end{tabular}}
\end{center}
Setiap kelas terdiri atas semua matriks $\nbyn{2}$ dengan rank yang sama.
Hanya terdapat satu matriks berank~nol.
Kedua kelas lainnya memiliki tak hingga banyak anggota; kita hanya
memperlihatkan wakil kanoniknya.
\end{example}

Salah satu keunggulan wakil dalam
\nearbytheorem{th:CanonFormForMatEquiv} adalah bahwa kita dapat memahami
pemetaan linear sepenuhnya ketika dinyatakan dengan cara berikut.
Jika basis-basisnya adalah
\( B=\sequence{\vec{\beta}_1,\dots,\vec{\beta}_n} \)
dan
\( D=\sequence{\vec{\delta}_1,\dots,\vec{\delta}_m} \)
maka tindakan pemetaannya adalah
\begin{equation*}
  c_1\vec{\beta}_1+\dots+c_k\vec{\beta}_k+c_{k+1}\vec{\beta}_{k+1}+\dots
     +c_n\vec{\beta}_n
  \;\mapsto\;
  c_1\vec{\delta}_1+\dots+c_k\vec{\delta}_k+\zero+\cdots+\zero
\end{equation*}
dengan \( k \) sebagai ranknya.
Jadi, kita dapat memandang sebarang pemetaan linear sebagai proyeksi.
\begin{equation*}
  \colvec{c_1 \\ \vdots \\ c_k \\ c_{k+1} \\ \vdots \\ c_n}_B
  \quad\longmapsto\quad
  \colvec{c_1 \\ \vdots \\ c_k \\ 0 \\ \vdots \\ 0}_D
\end{equation*}
% Of course, ``understanding'' a map expressed in this way 
% requires that we understand the relationship between \( B \) and \( D \).
% Nonetheless,
% this is a good classification of linear maps.

\begin{exercises}
  \recommended \item 
    Tentukan apakah pasangan-pasangan berikut ekuivalen secara matriks.
    \begin{exparts}
      \partsitem \( \begin{mat}[r]
                 1  &3  &0  \\
                 2  &3  &0
               \end{mat} \),
            \( \begin{mat}[r]
                 2  &2  &1  \\
                 0  &5  &-1
               \end{mat} \)
      \partsitem \( \begin{mat}[r]
                 0  &3  \\
                 1  &1
               \end{mat} \),
            \( \begin{mat}[r]
                 4  &0  \\
                 0  &5
               \end{mat} \)
      \partsitem \( \begin{mat}[r]
                 1  &3  \\
                 2  &6
               \end{mat} \),
            \( \begin{mat}[r]
                 1  &3  \\
                 2  &-6
               \end{mat} \)
    \end{exparts}
    \begin{answer}
      \begin{exparts}
        \partsitem Ya, masing-masing memiliki rank dua.
        \partsitem Ya, keduanya memiliki rank yang sama.
        \partsitem Tidak, keduanya memiliki rank berbeda.
      \end{exparts}  
    \end{answer}
  \item Manakah di antara matriks-matriks berikut yang ekuivalen secara matriks
    satu sama lain?
    \begin{exparts*}
      \partsitem 
         $\begin{mat}
            1 &2 &3 \\
            4 &5 &6 \\
            7 &8 &9
         \end{mat}$
      \partsitem 
         $\begin{mat}
           1 &3 \\
          -1 &-3
         \end{mat}$
      \partsitem 
         $\begin{mat}
            -5 &1 &0 \\
            -1 &0 &1
         \end{mat}$
      \partsitem 
         $\begin{mat}
            0 &-1 \\
            0 &5
         \end{mat}$
      \partsitem 
         $\begin{mat}
            1 &0 &1 \\
            2 &0 &2 \\
            1 &3 &1
         \end{mat}$
      \partsitem 
         $\begin{mat}
            3  &1   &0 \\
            9  &3   &0 \\
            -3 &-1  &0
         \end{mat}$
    \end{exparts*}
    \begin{answer}
      Kelompokkan matriks-matriks itu ke dalam kelas-kelas yang dicirikan oleh
      syarat bahwa semua matriks dalam kelas yang sama memiliki ukuran dan rank
      yang sama.
      \begin{exparts}
        \partsitem $\nbyn{3}$, rank~$2$
        \partsitem $\nbyn{2}$, rank~$1$
        \partsitem $\nbym{2}{3}$, rank~$2$
        \partsitem $\nbyn{3}$, rank~$2$
        \partsitem $\nbyn{3}$, rank~$1$
      \end{exparts}
    \end{answer}
  \recommended \item
    Temukan wakil kanonik kelas ekuivalensi matriks bagi setiap matriks.
    \begin{exparts*}
      \partsitem \( \begin{mat}[r]
                 2  &1  &0  \\
                 4  &2  &0
               \end{mat} \)
      \partsitem \( \begin{mat}[r]
                 0  &1  &0  &2  \\
                 1  &1  &0  &4  \\
                 3  &3  &3  &-1
               \end{mat} \)
    \end{exparts*}
    \begin{answer}
      Kita hanya perlu menghitung rank masing-masing.
      \begin{exparts*}
        \partsitem \( \begin{mat}[r]
                   1  &0  &0  \\
                   0  &0  &0
                 \end{mat} \)
        \partsitem \( \begin{mat}[r]
                   1  &0  &0  &0  \\
                   0  &1  &0  &0  \\
                   0  &0  &1  &0
                 \end{mat} \)
      \end{exparts*}  
    \end{answer}
  \item 
    Andaikan bahwa terhadap
    \begin{equation*}
      B=\stdbasis_2
      \qquad
      D=\sequence{\colvec[r]{1 \\ 1},\colvec[r]{1 \\ -1}}
    \end{equation*}
    transformasi \( \map{t}{\Re^2}{\Re^2} \) direpresentasikan oleh matriks
    berikut.
    \begin{equation*}
      \begin{mat}[r]
        1  &2  \\
        3  &4
      \end{mat}
    \end{equation*}
    Gunakan matriks-matriks perubahan basis untuk merepresentasikan \( t \)
    terhadap setiap pasangan berikut.
    \begin{exparts}
      \partsitem \( \hat{B}=\sequence{\colvec[r]{0 \\ 1},\colvec[r]{1 \\ 1}} \),
        \( \hat{D}=\sequence{\colvec[r]{-1 \\ 0},\colvec[r]{2 \\ 1}} \)
      \partsitem \( \hat{B}=\sequence{\colvec[r]{1 \\ 2},\colvec[r]{1 \\ 0}} \),
        \( \hat{D}=\sequence{\colvec[r]{1 \\ 2},\colvec[r]{2 \\ 1}} \)
    \end{exparts}
    \begin{answer}
      Ingat kembali diagram dan rumus berikut.
      \begin{equation*}
        \begin{CD}
          \Re^2_{\wrt{B}}                   @>t>T>        \Re^2_{\wrt{D}}       \\
          @V\text{\scriptsize$\identity$} VV   @V\text{\scriptsize$\identity$} VV \\
          \Re^2_{\wrt{\hat{B}}}             @>t>\hat{T}>  \Re^2_{\wrt{\hat{D}}}
        \end{CD}
        \qquad \hat{T}=
         \rep{\identity}{D,\hat{D}}\cdot T\cdot \rep{\identity}{\hat{B},B}
      \end{equation*}
      \begin{exparts}
        \partsitem Kedua persamaan berikut
          \begin{equation*}
            \colvec[r]{1 \\ 1}=1\cdot\colvec[r]{-1 \\ 0}
                            +1\cdot\colvec[r]{2 \\ 1}
            \qquad
            \colvec[r]{1 \\ -1}=(-3)\cdot\colvec[r]{-1 \\ 0}
                            +(-1)\cdot\colvec[r]{2 \\ 1}
          \end{equation*}
          menunjukkan bahwa
          \begin{equation*}
            \rep{\identity}{D,\hat{D}}=\begin{mat}[r]
              1  &-3  \\
              1  &-1
            \end{mat}
          \end{equation*}
          dan kedua persamaan berikut juga
          \begin{equation*}
            \colvec[r]{0 \\ 1}=0\cdot\colvec[r]{1 \\ 0}
                            +1\cdot\colvec[r]{0 \\ 1}
            \qquad
            \colvec[r]{1 \\  1}=1\cdot\colvec[r]{1 \\ 0}
                            +1\cdot\colvec[r]{0 \\ 1}
          \end{equation*}
          memberikan matriks nonsingular lainnya.
          \begin{equation*}
            \rep{\identity}{\hat{B},B}=\begin{mat}[r]
              0  &1  \\
              1  &1
            \end{mat}
          \end{equation*}
          Dengan demikian, jawabannya adalah
          \begin{equation*}
            \hat{T}=
            \begin{mat}[r]
              1  &-3  \\
              1  &-1
            \end{mat}
            \begin{mat}[r]
              1  &2  \\
              3  &4
            \end{mat}
            \begin{mat}[r]
              0  &1  \\
              1  &1
            \end{mat}
            =\begin{mat}[r]
              -10  &-18  \\
              -2   &-4 
            \end{mat}
          \end{equation*}
          Meskipun tidak mutlak diperlukan, pemeriksaan dapat meyakinkan kita.
          Tetapkan secara sebarang
          \begin{equation*}
            \vec{v}=\colvec[r]{3 \\ 2}
          \end{equation*}
kita memperoleh
          \begin{equation*}
            \rep{\vec{v}}{B}=\colvec[r]{3 \\ 2}_B
            \qquad
            \begin{mat}[r]
              1  &2  \\
              3  &4
            \end{mat}_{B,D}
            \colvec[r]{3 \\ 2}_B
            =\colvec[r]{7 \\ 17}_D
          \end{equation*}
          sehingga $t(\vec{v})$ adalah
          \begin{equation*}
            7\cdot\colvec[r]{1 \\ 1}+17\cdot\colvec[r]{1 \\ -1}=\colvec[r]{24 \\ -10}
          \end{equation*}
          Perhitungan terhadap $\hat{B},\hat{D}$ dimulai dengan
          \begin{equation*}
            \rep{\vec{v}}{\hat{B}}=\colvec[r]{-1 \\ 3}_{\hat{B}}
            \qquad
            \begin{mat}[r]
              -10  &-18  \\
               -2  &-4
            \end{mat}_{\hat{B},\hat{D}}
            \colvec[r]{-1 \\ 3}_{\hat{B}}
            =\colvec[r]{-44 \\ -10}_{\hat{D}}
          \end{equation*}
          lalu pemeriksaan berikut menunjukkan bahwa hasilnya sama.
          \begin{equation*}
            -44\cdot\colvec[r]{-1 \\ 0}-10\cdot\colvec[r]{2 \\ 1}=\colvec[r]{24 \\ -10}
          \end{equation*}
    \partsitem Kedua persamaan berikut
      \begin{equation*}
        \colvec[r]{1 \\ 1}=\frac{1}{3}\cdot\colvec[r]{1 \\ 2}
                  +\frac{1}{3}\cdot\colvec[r]{2 \\ 1}
        \qquad
        \colvec[r]{1 \\ -1}=-1\cdot\colvec[r]{1 \\ 2}
                  +1\cdot\colvec[r]{2 \\ 1}
      \end{equation*}
      menunjukkan bahwa
      \begin{equation*}
        \rep{\identity}{D,\hat{D}}=\begin{mat}[r]
          1/3  &-1  \\
          1/3  &1
        \end{mat}
      \end{equation*}
      dan kedua persamaan berikut
      \begin{equation*}
        \colvec[r]{1 \\ 2}=1\cdot\colvec[r]{1 \\ 0}
                  +2\cdot\colvec[r]{0 \\ 1}
        \qquad
        \colvec[r]{1 \\ 0}=1\cdot\colvec[r]{1 \\ 0}
                  +0\cdot\colvec[r]{0 \\ 1}
      \end{equation*}
      menunjukkan hasil berikut.
      \begin{equation*}
        \rep{\identity}{\hat{B},B}=\begin{mat}[r]
          1  &1  \\
          2  &0
        \end{mat}
      \end{equation*}
      Dengan kedua matriks tersebut, perubahannya berlangsung sebagai berikut.
      \begin{equation*}
        \hat{T}=\begin{mat}[r]
          1/3  &-1  \\
          1/3  &1
        \end{mat}
        \begin{mat}[r]
          1  &2  \\
          3  &4
        \end{mat}
        \begin{mat}[r]
          1  &1  \\
          2  &0
        \end{mat}
        =\begin{mat}[r]
          -28/3  &-8/3  \\
          38/3   &10/3
        \end{mat}
      \end{equation*}
      Seperti pada butir sebelumnya, pemeriksaan membantu memastikan bahwa
      perhitungan ini bebas dari kesalahan.
      Sebagai contoh, kita dapat menetapkan vektor
      \begin{equation*}
        \vec{v}=\colvec[r]{-1 \\ 2}
      \end{equation*}
      (vektor ini dipilih secara sebarang).
      Sekarang kita memiliki
      \begin{equation*}
        \rep{\vec{v}}{B}=\colvec[r]{-1 \\ 2}
        \qquad
        \begin{mat}[r]
          1  &2  \\
          3  &4
        \end{mat}_{B,D}
        \colvec[r]{-1  \\ 2}_{B}
        =\colvec[r]{3  \\ 5}_D
      \end{equation*}
      sehingga $t(\vec{v})$ adalah vektor berikut.
      \begin{equation*}
        3\cdot\colvec[r]{1 \\ 1}+5\cdot\colvec[r]{1 \\ -1}=\colvec[r]{8 \\ -2}
      \end{equation*}
      Terhadap $\hat{B},\hat{D}$, pertama-tama kita menghitung
      \begin{equation*}
        \rep{\vec{v}}{\hat{B}}=\colvec[r]{1 \\ -2}
        \qquad
        \begin{mat}[r]
          -28/3  &-8/3  \\
          38/3   &10/3
        \end{mat}_{\hat{B},\hat{D}}
        \colvec[r]{1 \\ -2}_{\hat{B}}
        =\colvec[r]{-4 \\ 6}_{\hat{D}}
      \end{equation*}
      dan, sebagaimana diharapkan, hasil $t(\vec{v})$ tetap sama.
      \begin{equation*}
        -4\cdot\colvec[r]{1 \\ 2}+6\cdot\colvec[r]{2 \\ 1}=\colvec[r]{8 \\ -2}
      \end{equation*}
      \end{exparts} 
    \end{answer}
  \item 
    Berapakah ukuran \( P \) dan \( Q \) dalam persamaan $\hat{H}=PHQ$?
    \begin{answer}
      Jika \( H \) dan \( \hat{H} \) berukuran \( \nbym{m}{n} \), maka
      matriks \( P \) berukuran \( \nbyn{m} \), sedangkan \( Q \) berukuran
      \( \nbyn{n} \).
    \end{answer}
  \recommended \item
    Tinjau ruang $V=\polyspace_2$ dan~$W=\matspace_{\nbyn{2}}$ dengan
    basis-basis berikut.
    \begin{gather*}
      B=\sequence{1, 1+x, 1+x^2}
      \qquad
      D=\sequence{
        \begin{mat}
          0 &0 \\
          0 &1
        \end{mat},
        \begin{mat}
          0 &0 \\
          1 &1
        \end{mat},
        \begin{mat}
          0 &1 \\
          1 &1
        \end{mat},
        \begin{mat}
          1 &1 \\
          1 &1
        \end{mat}
         }                        \\
      \hat{B}=\sequence{1, x, x^2}
      \qquad
      \hat{D}=\sequence{
        \begin{mat}
          -1 &0 \\
          0 &0
        \end{mat},
        \begin{mat}
          0 &-1 \\
          0 &0
        \end{mat},
        \begin{mat}
          0 &0 \\
          1 &0
        \end{mat},
        \begin{mat}
          0 &0 \\
          0 &1
        \end{mat}
         }
    \end{gather*}
    Kita akan menemukan $P$ dan~$Q$ untuk mengubah representasi suatu pemetaan
    terhadap $B,D$ menjadi representasi terhadap $\hat{B},\hat{D}$.
    \begin{exparts}
      \partsitem Gambarkan diagram panah yang sesuai.
      \partsitem Hitung $P$ dan~$Q$.
    \end{exparts}
    \begin{answer}
      \begin{exparts}
        \partsitem
          Berikut diagram panahnya.
          \begin{equation*}
            \begin{CD}
              V_{\wrt{B}}                   @>h>H>        W_{\wrt{D}}       \\
              @V{\text{\scriptsize$\identity$}} VV      @V{\text{\scriptsize$\identity$}} VV \\
              V_{\wrt{\hat{B}}}             @>h>\hat{H}>   W_{\wrt{\hat{D}}}
            \end{CD}
          \end{equation*}
          Diagram tersebut memberikan $\hat{H}=PHQ$, dengan
           $Q=\rep{\identity}{\hat{B},B}$ dan $P=\rep{\identity}{D,\hat{D}}$
            (ingat bahwa operasi yang dilakukan lebih dahulu ditulis di sebelah
            kanan).
        \partsitem
          Kita menginginkan $Q=\rep{\identity}{\hat{B},B}$ dan
          $P=\rep{\identity}{D,\hat{D}}$.
          Untuk $Q$, kita melakukan perhitungan berikut (di sini dilakukan
          secara langsung).
          \begin{equation*}
            \rep{\identity(1)}{B}=\colvec{1 \\ 0 \\ 0}
            \quad
            \rep{\identity(x)}{B}=\colvec{-1 \\ 1 \\ 0}
            \quad
            \rep{\identity(x^2)}{B}=\colvec{-1 \\ 0 \\ 1}
          \end{equation*}
          Perhitungan berikut memberikan $P$.
          \begin{multline*}
            \rep{\identity(
              \begin{mat}
                0 &0 \\
                0 &1
              \end{mat})}{\hat{D}}=\colvec{0 \\ 0 \\ 0 \\ 1} 
            \quad
            \rep{\identity(    
              \begin{mat}
                0 &0 \\
                1 &1
              \end{mat})}{\hat{D}}=\colvec{0 \\ 0 \\ 1 \\ 1} 
            \\
            \rep{\identity(    
              \begin{mat}
                0 &1 \\
                1 &1
              \end{mat})}{\hat{D}}=\colvec{0 \\ -1 \\ 1 \\ 1}
            \quad
            \rep{\identity(    
              \begin{mat}
                1 &1 \\
                1 &1
              \end{mat})}{\hat{D}}=\colvec{-1 \\ -1 \\ 1 \\ 1} 
          \end{multline*}
          Jadi, jawabannya adalah
          \begin{equation*}
            P=
            \begin{mat}
              0 &0 &0  &-1 \\
              0 &0 &-1 &-1  \\
              0 &1 &1  &1  \\
              1 &1 &1  &1
            \end{mat}
            \qquad
            Q=\begin{mat}
              1 &-1 &-1 \\
              0 &1  &0  \\
              0 &0  &1
            \end{mat}
          \end{equation*}      
      \end{exparts}
    \end{answer}
  \recommended \item Temukan matriks perubahan basis $Q$ dan~$P$ yang akan
   mengubah representasi suatu $\map{t}{\Re^2}{\Re^2}$ terhadap $B,D$
   menjadi representasi terhadap $\hat{B},\hat{D}$.
   \begin{equation*}
     B=\sequence{\colvec{1 \\ 0}, \colvec{1 \\ 1}}\quad
     D=\sequence{\colvec{0 \\ 1}, \colvec{-1 \\ 0}}\quad
     \hat{B}=\stdbasis_2\quad
     \hat{D}=\sequence{\colvec{1 \\ -1}, \colvec{0 \\ 1}}
   \end{equation*}
    \begin{answer}
          Berikut diagram panahnya.
          \begin{equation*}
            \begin{CD}
              V_{\wrt{B}}                   @>t>T>        W_{\wrt{D}}       \\
              @V{\text{\scriptsize$\identity$}} VV      @V{\text{\scriptsize$\identity$}} VV \\
              V_{\wrt{\hat{B}}}             @>t>\hat{T}>   W_{\wrt{\hat{D}}}
            \end{CD}
          \end{equation*}
          Dengan $Q=\rep{\identity}{\hat{B},B}$ dan
          $P=\rep{\identity}{D,\hat{D}}$
          persamaannya adalah $\hat{T}=PTQ$.
 
          Berikut perhitungan untuk $Q=\rep{\identity}{\hat{B},B}$
          (dilakukan secara langsung).
          \begin{equation*}
            \rep{\identity(\colvec{1 \\ 0})}{B}=\colvec{1 \\ 0}
            \quad
            \rep{\identity(\colvec{0 \\ 1})}{B}=\colvec{-1 \\ 1}
          \end{equation*}
          Jadi, kita memperoleh
          \begin{equation*}
            Q=\begin{mat}
              1 &-1 \\
              0 &1  
             \end{mat}
          \end{equation*}
          Perhitungan berikut memberikan
          $P=\rep{\identity}{D,\hat{D}}$.
          \begin{equation*}
            \rep{\identity(\colvec{0 \\ 1})}{\hat{D}}=\colvec{0 \\ 1}
            \quad
            \rep{\identity(\colvec{-1 \\ 0})}{\hat{D}}=\colvec{-1 \\ -1}
          \end{equation*}
          Susun keduanya sebagai kolom-kolom matriks lainnya.
          \begin{equation*}
            P=\begin{mat}
              0 &-1 \\
              1 &-1  
             \end{mat}
          \end{equation*}
    \end{answer}
  \recommended \item Temukan $P$ dan~$Q$ untuk menyatakan $H$ melalui $PHQ$
    sebagai matriks blok identitas parsial.
    \begin{equation*}
      H=
      \begin{mat}
        2 &1  &1  \\
        3 &-1 &0  \\
        1 &3  &2 
      \end{mat}
    \end{equation*}
    \begin{answer}
      Metode Gauss menghasilkan
      \begin{equation*}
          \begin{mat}
            2 &1  &1  \\
            3 &-1 &0  \\
            1 &3  &2 
          \end{mat}
          \grstep[-(1/2)\rho_1+\rho_3]{-(3/2)\rho_1+\rho_2}
          \grstep{\rho_2+\rho_3}
          \grstep[-(2/5)\rho_2]{(1/2)\rho_1}
          \begin{mat}
            1 &1/2  &1/2  \\
            0 &1    &3/5  \\
            0 &0    &0 
          \end{mat}
      \end{equation*}
      Operasi kolom menyelesaikan proses mencapai bentuk kanonik bagi
      ekuivalensi matriks.
      \begin{equation*}
        \grstep{-(3/5)\text{kol}_2+\text{kol}_3}
        \grstep[-(1/5)\text{kol}_1+\text{kol}_3]{-(1/2)\text{kol}_1+\text{kol}_2}
          \begin{mat}
            1 &0    &0  \\
            0 &1 &0  \\
            0 &0  &0 
          \end{mat}
      \end{equation*}
      Dengan demikian, kedua matriksnya adalah
      \begin{align*}
        P&=
        \begin{mat}
          1 &0 &0 \\
          0 &-2/5 &0 \\
          0 &0 &1
         \end{mat}
        \begin{mat}
          1/2 &0 &0 \\
          0 &1 &0 \\
          0 &0 &1
         \end{mat}
        \begin{mat}
          1 &0 &0 \\
          0 &1 &0 \\
          0 &1 &1
         \end{mat}
        \begin{mat}
          1    &0 &0 \\
          0    &1 &0 \\
          -1/2 &0 &1
         \end{mat}
        \begin{mat}
          1    &0 &0 \\
          -3/2 &1 &0 \\
          0    &0 &1
         \end{mat}                                 \\ 
         &=
         \begin{mat}
           1/2   &0     &0 \\
           3/5   &-2/5  &0 \\
           -2    &1     &1
         \end{mat}                              \\
        Q&=
        \begin{mat}
          1 &0 &0 \\
          0 &1 &-3/5 \\
          0 &0 &1
        \end{mat}
        \begin{mat}
          1 &-1/2 &0 \\
          0 &1    &0 \\
          0 &0    &1
        \end{mat}
        \begin{mat}
          1 &0 &-1/5 \\
          0 &1 &0 \\
          0 &0 &1
        \end{mat}
        =
        \begin{mat}
          1 &-1/2 &-1/5 \\
          0 &1    &-3/5 \\
          0 &0    &1
        \end{mat}
      \end{align*}      
    \end{answer}
  \recommended \item
    Gunakan \nearbytheorem{th:CanonFormForMatEquiv} untuk menunjukkan bahwa
    suatu matriks persegi nonsingular jika dan hanya jika ekuivalen dengan
    matriks identitas.
    \begin{answer}
        Sebarang matriks \( \nbyn{n} \) nonsingular jika dan hanya jika
        memiliki rank \( n \); menurut
        \nearbytheorem{th:CanonFormForMatEquiv}, hal ini berlaku jika dan hanya
        jika matriks tersebut ekuivalen secara matriks dengan matriks
        $\nbyn{n}$ yang semua entri diagonalnya bernilai satu.
    \end{answer}
  \item
    Tunjukkan bahwa jika \( A \) merupakan matriks persegi nonsingular dan
    \( P \) serta \( Q \) merupakan matriks-matriks persegi nonsingular
    sedemikian sehingga \( PAQ=I \), maka \( QP=A^{-1} \).
    \begin{answer}
      Jika \( PAQ=I \), maka \( QPAQ=Q \), sehingga \( QPA=I \), dan karena
      itu \( QP=A^{-1} \).
    \end{answer}
  \item
    Mengapa \nearbytheorem{th:CanonFormForMatEquiv} tidak menunjukkan bahwa
    setiap matriks dapat didiagonalkan (lihat
    \nearbyexample{ex:DiagizedMat})?
    \begin{answer}
      Menurut definisi setelah \nearbyexample{ex:DiagizedMat}, suatu matriks
      $M=\rep{t}{B,D}$ dapat didiagonalkan jika matriks tersebut
      merepresentasikan transformasi dengan sifat bahwa terdapat basis
      $\hat{B}$ sedemikian sehingga $\rep{t}{\hat{B},\hat{B}}$ merupakan
      matriks diagonal\Dash basis awal dan akhirnya harus sama.
      Namun, \nearbytheorem{th:CanonFormForMatEquiv} hanya menyatakan bahwa
      terdapat $\hat{B}$ dan $\hat{D}$ sedemikian sehingga kita dapat beralih
      ke representasi $\rep{t}{\hat{B},\hat{D}}$ dan memperoleh matriks
      diagonal. Tidak ada alasan untuk menganggap bahwa kita dapat memilih
      $\hat{B}$ dan $\hat{D}$ agar keduanya sama.
    \end{answer}
  \item 
    Apakah transpos dari matriks-matriks yang ekuivalen secara matriks juga
    harus ekuivalen secara matriks?
    \begin{answer}
      Ya. Rank baris sama dengan rank kolom, sehingga rank transpos sama dengan
      rank matriksnya. Matriks-matriks berukuran sama dengan rank yang sama
      ekuivalen secara matriks.
    \end{answer}
  \item 
    Apa yang terjadi dalam \nearbytheorem{th:CanonFormForMatEquiv} jika
    \( k=0 \)?
    \begin{answer}
      Hanya matriks nol yang memiliki rank nol.
    \end{answer}
  \item \label{exer:MatEqIsEqRel}
    Tunjukkan bahwa ekuivalensi matriks merupakan relasi ekuivalensi.
    \begin{answer}
      Untuk refleksivitas, guna menunjukkan bahwa sebarang matriks ekuivalen
      secara matriks dengan dirinya sendiri, ambil \( P \) dan \( Q \) sebagai
      matriks identitas. Untuk simetri, jika \( H_1=PH_2Q \), maka
      \( H_2=P^{-1}H_1Q^{-1} \) (inversnya ada karena $P$ dan $Q$
      nonsingular). Terakhir, untuk transitivitas, andaikan
      \( H_1=P_2H_2Q_2 \) dan \( H_2=P_3H_3Q_3 \).
      Substitusi kemudian memberikan
      \( H_1=P_2(P_3H_3Q_3)Q_2=(P_2P_3)H_3(Q_3Q_2) \).
      Hasil kali matriks-matriks nonsingular bersifat nonsingular (kita telah
      menunjukkan bahwa hasil kali matriks-matriks yang dapat diinvers juga
      dapat diinvers;~bahkan kita telah menunjukkan cara menghitung inversnya),
      sehingga \( H_1 \) ekuivalen secara matriks dengan \( H_3 \).
     \end{answer}
  \recommended \item 
    Tunjukkan bahwa matriks nol merupakan satu-satunya anggota kelas
    ekuivalensi matriksnya. Adakah matriks lain dengan sifat tersebut?
    \begin{answer}
      Menurut \nearbytheorem{th:CanonFormForMatEquiv}, matriks nol merupakan
      satu-satunya anggota kelasnya karena merupakan satu-satunya matriks
      \( \nbym{m}{n} \) berank nol.
      Tidak ada matriks lain yang menjadi satu-satunya anggota kelasnya;
      sebarang kelipatan skalar tak nol dari suatu matriks memiliki rank yang
      sama dengan matriks tersebut.
    \end{answer}
  \item 
    Apa saja kelas ekuivalensi matriks bagi matriks transformasi pada
    \( \Re^1 \)?
    \( \Re^3 \)?
    \begin{answer}
      Terdapat dua kelas ekuivalensi matriks bagi matriks \( \nbyn{1} \)\Dash
      kelas matriks berank nol dan kelas matriks berank satu.
      Matriks-matriks \( \nbyn{3} \) terbagi ke dalam empat kelas ekuivalensi
      matriks.
    \end{answer}
  \item 
    Berapa banyak kelas ekuivalensi matriks yang ada?
    \begin{answer}
      Bagi matriks \( \nbym{m}{n} \), terdapat satu kelas untuk setiap rank
      yang mungkin:~jika \( k \) merupakan nilai minimum dari \( m \) dan
      \( n \), terdapat kelas bagi matriks berank \( 0 \), \( 1 \), \ldots,
      \( k \). Jadi, terdapat \( k+1 \) kelas.
      (Tentu saja, jika semua ukuran matriks digabungkan, kita memperoleh tak
      hingga banyak kelas.)
    \end{answer}
  \item 
    Apakah kelas-kelas ekuivalensi matriks tertutup terhadap perkalian skalar?
    Terhadap penjumlahan?
    \begin{answer}
      Kelas-kelas tersebut tertutup terhadap perkalian skalar tak nol karena
      kelipatan skalar tak nol dari suatu matriks memiliki rank yang sama
      dengan matriks tersebut. Kelas-kelas itu tidak tertutup terhadap
      penjumlahan; sebagai contoh, \( H+(-H) \) memiliki rank nol.
    \end{answer}
  \item 
    Misalkan \( \map{t}{\Re^n}{\Re^n} \) direpresentasikan oleh \( T \)
    terhadap \( \stdbasis_n,\stdbasis_n \).
    \begin{exparts}
      \partsitem Temukan $\rep{t}{B,B}$ dalam kasus khusus berikut.
        \begin{equation*}
          T=\begin{mat}[r]
            1  &1  \\
            3  &-1
          \end{mat}
          \qquad
          B=\sequence{\colvec[r]{1  \\ 2},
                      \colvec[r]{-1 \\ -1}}
        \end{equation*}
      \partsitem Deskripsikan $\rep{t}{B,B}$ dalam kasus umum dengan
        \( B=\basis{\beta}{n} \).
    \end{exparts}
    \begin{answer}
        Berikut diagramnya.
        \begin{equation*}
          \begin{CD}
            \Re^2_{\wrt{\stdbasis_2}}         @>t>T>     \Re^2_{\wrt{\stdbasis_2}}    \\
            @V\text{\scriptsize$\identity$} VV   @V\text{\scriptsize$\identity$} VV \\
           \Re^2_{\wrt{B}}        @>t>\hat{T}>  \Re^2_{\wrt{B}}
          \end{CD}
        \end{equation*}
        Terdapat dua cara untuk bergerak dari kiri bawah ke kanan bawah.
        Cara pertama adalah secara langsung dengan menggunakan
        $\hat{T}=\rep{t}{B,B}$.
        Cara kedua bergerak naik, lalu mendatar, kemudian turun, dengan
        $\rep{\identity}{\stdbasis_2,B}
                          \cdot T
                          \cdot \rep{\identity}{B,\stdbasis_2}$
        (ingat bahwa matriks-matriks tersebut ditulis dari kanan ke kiri,
        sehingga matriks untuk gerak ``naik'' berada di sebelah kanan; dengan
        demikian, ketika diterapkan pada $\rep{\vec{v}}{B}$, matriks yang
        diterapkan pertama kali adalah $\rep{\identity}{B,\stdbasis_2}$).
        Kita menulis $S$ untuk $\rep{\identity}{B,\stdbasis_2}$; di antara
        kedua matriks perubahan basis, kita memilih matriks ini karena lebih
        mudah dihitung. Jadi, kita memiliki $\hat{T}=S^{-1}TS$.
      \begin{exparts}
        \partsitem Kita memiliki
          \begin{equation*}
            \rep{\identity}{B,\stdbasis_2}=S=
            \begin{mat}[r]
              1  &-1  \\
              2  &-1
            \end{mat}
          \end{equation*}
          dan
          \begin{equation*}
            \rep{\identity}{\stdbasis_2,B}
            =\rep{\identity}{B,\stdbasis_2}^{-1}=S^{-1}=
            \begin{mat}[r]
              1  &-1  \\
              2  &-1
            \end{mat}^{-1}
            =\begin{mat}[r]
              -1  &1  \\
              -2  &1 
            \end{mat}
          \end{equation*}
          sehingga jawabannya adalah
          \begin{equation*}
            \rep{t}{B,B}
            =
            \begin{mat}[r]
              -1  &1  \\
              -2  &1 
            \end{mat}
            \begin{mat}[r]
              1  &1  \\
              3  &-1   
            \end{mat}
            \begin{mat}[r]
              1  &-1  \\
              2  &-1
            \end{mat}
            =\begin{mat}[r]
              -2  &0   \\
              -5  &2 
            \end{mat}
          \end{equation*}
          Sebagai pemeriksaan singkat, kita dapat mengambil vektor secara acak
          \begin{equation*}
            \vec{v}=\colvec[r]{4  \\  5}
          \end{equation*}
          yang memberikan
          \begin{equation*}
            \rep{\vec{v}}{\stdbasis_2}=\colvec[r]{4 \\ 5}
            \qquad
            \begin{mat}[r]
              1  &1  \\
              3  &-1   
            \end{mat}
            \colvec[r]{4 \\ 5}
            =\colvec[r]{9 \\ 7}=t(\vec{v})
          \end{equation*}
          sedangkan perhitungan terhadap $B,B$
          \begin{equation*}
            \rep{\vec{v}}{B}=\colvec[r]{1 \\ -3}
            \qquad
            \begin{mat}[r]
              -2  &0   \\
              -5  &2 
            \end{mat}_{B,B}
            \colvec[r]{1 \\ -3}_B
            =\colvec[r]{-2 \\ -11}_B
          \end{equation*}
          menghasilkan jawaban yang sama.
          \begin{equation*}
            -2\cdot\colvec[r]{1 \\ 2}-11\cdot\colvec[r]{-1 \\ -1}
               =\colvec[r]{9 \\ 7}
          \end{equation*}
       \partsitem Seperti pada butir pertama pertanyaan ini,
          \begin{equation*}
            S=\rep{\identity}{B,\stdbasis_2}
            =\begin{pmat}{c|c|c}
              \vec{\beta}_1 &\;\cdots\; &\vec{\beta}_n
            \end{pmat}
            \qquad
            \rep{\identity}{\stdbasis_2,B}=S^{-1}=
             \rep{\identity}{B,\stdbasis_2}^{-1}
          \end{equation*}
          sehingga, dengan menulis $S$ untuk matriks yang kolom-kolomnya
          merupakan vektor-vektor basis, kita memiliki
          $\rep{t}{B,B}=\hat{T}=S^{-1}TS$.
      \end{exparts}
    \end{answer}
  \item 
    \begin{exparts}
       \partsitem Misalkan \( V \) memiliki basis \( B_1 \) dan \( B_2 \),
         sedangkan \( W \) memiliki basis \( D \).
         Untuk \( \map{h}{V}{W} \), temukan rumus yang menghitung
         \( \rep{h}{B_2,D} \) dari \( \rep{h}{B_1,D} \).
       \partsitem Ulangi pertanyaan sebelumnya dengan satu basis bagi \( V \)
         dan dua basis bagi \( W \).
     \end{exparts}
     \begin{answer}
       \begin{exparts}
         \partsitem Bentuk diagram panah yang disesuaikan adalah
           \begin{equation*}
             \begin{CD}
               V_{\wrt{B_1}}                   @>h>H>        W_{\wrt{D}}       \\
               @V\text{\scriptsize$\identity$} VQV      @V\text{\scriptsize$\identity$} VPV \\
               V_{\wrt{B_2}}             @>h>\hat{H}>  W_{\wrt{D}}
             \end{CD}
           \end{equation*}
           Karena tidak perlu mengubah basis dalam \( W \) (atau kita dapat
           mengatakan bahwa matriks perubahan basis $P$ merupakan identitas),
           kita memiliki
           \( \rep{h}{B_2,D}=\rep{h}{B_1,D}\cdot Q \), dengan
           \( Q=\rep{\identity}{B_2,B_1} \).
         \partsitem Dalam kasus ini, diagram panahnya adalah
           \begin{equation*}
             \begin{CD}
               V_{\wrt{B}}                   @>h>H>   W_{\wrt{D_1}}       \\
               @V\text{\scriptsize$\identity$} VQV  @V\text{\scriptsize$\identity$} VPV \\
               V_{\wrt{B}}             @>h>\hat{H}>  W_{\wrt{D_2}}
             \end{CD}
           \end{equation*}
           Kita memiliki \( \rep{h}{B,D_2}=P\cdot \rep{h}{B,D_1} \), dengan
           \( P=\rep{\identity}{D_1,D_2} \).
       \end{exparts}  
      \end{answer}
  \item 
    \begin{exparts}
      \partsitem Jika dua matriks ekuivalen secara matriks dan dapat diinvers,
        apakah invers keduanya harus ekuivalen secara matriks?
      \partsitem Jika dua matriks memiliki invers yang ekuivalen secara matriks,
        apakah kedua matriks tersebut harus ekuivalen secara matriks?
      \partsitem Jika dua matriks persegi ekuivalen secara matriks, apakah
        kuadrat keduanya harus ekuivalen secara matriks?
      \partsitem Jika dua matriks persegi memiliki kuadrat yang ekuivalen secara
        matriks, apakah kedua matriks itu harus ekuivalen secara matriks?
    \end{exparts}
    \begin{answer}
      \begin{exparts}
        \partsitem Berikut diagram panahnya beserta versi diagram tersebut bagi
          fungsi invers.
          \begin{equation*}
           \begin{CD}
             V_{\wrt{B}}                   @>h>H>      W_{\wrt{D}}       \\
             @V\text{\scriptsize$\identity$} VQV  @V\text{\scriptsize$\identity$} VPV \\
             V_{\wrt{\hat{B}}}             @>h>\hat{H}> W_{\wrt{\hat{D}}}
            \end{CD}
            \hspace*{4em}\qquad
           \begin{CD}
             V_{\wrt{B}}              @<h^{-1}<H^{-1}<    W_{\wrt{D}}       \\
             @V\text{\scriptsize$\identity$} VQV  @V\text{\scriptsize$\identity$} VPV \\
             V_{\wrt{\hat{B}}}        @<h^{-1}<\hat{H}^{-1}< W_{\wrt{\hat{D}}}
            \end{CD}
           \end{equation*}
           Ya, invers matriks-matriks tersebut merepresentasikan invers
           pemetaan-pemetaannya. Artinya, kita dapat bergerak dari kanan bawah
           ke kiri bawah dengan bergerak naik, lalu ke kiri, kemudian turun.
           Dengan kata lain, jika \( \hat{H}=PHQ \) (dengan \( P,Q \) dapat
           diinvers) serta \( H,\hat{H} \) dapat diinvers, maka
           \( \hat{H}^{-1}=Q^{-1}H^{-1}P^{-1} \).
        \partsitem Ya; ini merupakan bagian sebelumnya yang dinyatakan dengan
          istilah berbeda.
        \partsitem Tidak, kita memerlukan asumsi tambahan:~jika \( H \)
          merepresentasikan \( h \) terhadap basis awal dan akhir yang sama,
          yaitu \( B,B \), bagi suatu \( B \), maka \( H^2 \)
          merepresentasikan \( \composed{h}{h} \).
          Sebagai contoh khusus, kedua matriks berikut sama-sama berank satu
          sehingga ekuivalen secara matriks
          \begin{equation*}
             \begin{mat}[r]
               1  &0  \\
               0  &0
             \end{mat}
             \qquad
             \begin{mat}[r]
               0  &0  \\
               1  &0
             \end{mat}
          \end{equation*}
          tetapi kuadrat keduanya tidak ekuivalen secara matriks\Dash kuadrat
          matriks pertama berank satu, sedangkan kuadrat matriks kedua berank
          nol.
        \partsitem Tidak. Kedua matriks berikut tidak ekuivalen secara matriks,
          tetapi memiliki kuadrat yang ekuivalen secara matriks.
          \begin{equation*}
             \begin{mat}[r]
               0  &0  \\
               0  &0
             \end{mat}
             \qquad
             \begin{mat}[r]
               0  &0  \\
               1  &0
             \end{mat} 
          \end{equation*}
      \end{exparts}  
    \end{answer}
  \item
    Matriks-matriks persegi disebut \definend{serupa} jika merepresentasikan
    transformasi yang sama, tetapi masing-masing terhadap basis akhir dan awal
    yang sama. Artinya, \( \rep{t}{B_1,B_1} \) serupa dengan
    \( \rep{t}{B_2,B_2} \).
    \begin{exparts}
      \partsitem Berikan definisi keserupaan matriks seperti definisi dalam
        \nearbydefinition{def:MatEquiv}.
      \partsitem Buktikan bahwa matriks-matriks serupa ekuivalen secara matriks.
      \partsitem Tunjukkan bahwa keserupaan merupakan relasi ekuivalensi.
      \partsitem Tunjukkan bahwa jika \( T \) serupa dengan \( \hat{T} \),
        maka \( T^2 \) serupa dengan \( \hat{T}^2 \), kubiknya serupa, dan
        seterusnya.
        % \textit{Contrast with the prior exercise.}
      \partsitem Buktikan bahwa terdapat matriks-matriks yang ekuivalen secara
        matriks tetapi tidak serupa.
    \end{exparts}
    \begin{answer}
      \begin{exparts}
        \partsitem Diagram panah berikut menyarankan definisinya.
          \begin{equation*}
            \begin{CD}
              V_{\wrt{B_1}}                   @>t>T>        V_{\wrt{B_1}}       \\
              @V\text{\scriptsize$\identity$} VV  @V\text{\scriptsize$\identity$} VV \\
              V_{\wrt{B_2}}             @>t>\hat{T}>  V_{\wrt{B_2}}
            \end{CD}
          \end{equation*}
          Matriks \( T, \hat{T} \) disebut \definend{serupa} jika terdapat
          matriks nonsingular \( P \) sedemikian sehingga
          \( \hat{T}=P^{-1}TP \).
        \partsitem Dalam definisi ekuivalensi matriks, ambil pengali kiri
          sebagai \( P^{-1} \) dan pengali kanan sebagai \( P \).
        \partsitem \textit{Argumennya sama seperti dalam
          \nearbyexercise{exer:MatEqIsEqRel}.}
          Refleksivitas jelas: \( T=I^{-1}TI \).
          Simetri juga mudah: \( \hat{T}=P^{-1}TP \) menyiratkan
          \( T=P\hat{T}P^{-1} \) (kalikan persamaan pertama dari kanan dengan
          $P^{-1}$ dan dari kiri dengan $P$).
          Untuk transitivitas, andaikan \( T_1={P_2}^{-1}T_2P_2 \) dan
          \( T_2={P_3}^{-1}T_3P_3 \).
          Maka \( T_1={P_2}^{-1}({P_3}^{-1}T_3P_3)P_2
                     =({P_2}^{-1}{P_3}^{-1})T_3(P_3P_2) \), dan pembuktian
          selesai setelah memperhatikan bahwa \( P_3P_2 \) merupakan matriks
          yang dapat diinvers dengan invers \( {P_2}^{-1}{P_3}^{-1} \).
        \partsitem Andaikan \( \hat{T}=P^{-1}TP \). Untuk kuadrat,
          \( \hat{T}^2=(P^{-1}TP)(P^{-1}TP)
                      =P^{-1}T(PP^{-1})TP=P^{-1}T^2P \).
          Pangkat-pangkat yang lebih tinggi mengikuti melalui induksi.
        \partsitem Kedua matriks berikut ekuivalen secara matriks, tetapi
          kuadrat keduanya tidak ekuivalen secara matriks.
          \begin{equation*}
             \begin{mat}[r]
               1  &0  \\
               0  &0
             \end{mat}
             \qquad
             \begin{mat}[r]
               0  &0  \\
               1  &0
             \end{mat}
          \end{equation*}
          Menurut butir sebelumnya, keserupaan matriks dan ekuivalensi matriks
          dengan demikian merupakan dua hal berbeda.
      \end{exparts}  
   \end{answer}
\index{ekuivalensi matriks|)}
\index{perubahan basis|)}
\end{exercises}
